\documentclass[12pt]{article}
\usepackage[a4paper,margin=1in]{geometry}
\usepackage{iftex}
\ifXeTeX
\usepackage{fontspec}
\defaultfontfeatures{Ligatures=TeX}
\setmainfont{Times New Roman}
\else
\usepackage[T1]{fontenc}
\usepackage[utf8]{inputenc}
\usepackage{newtxtext}
\usepackage{newtxmath}
\fi
\usepackage{enumitem}
\usepackage{hyperref}
\usepackage{parskip}
\usepackage{microtype}
\usepackage[normalem]{ulem}
\usepackage{xcolor}

\newcommand{\practiceid}[1]{\texttt{#1}}
\newlength{\readinggutter}
\setlength{\readinggutter}{2.8em}
\newlength{\readingfirstindent}
\setlength{\readingfirstindent}{1.5em}
\newlength{\questionblockgap}
\setlength{\questionblockgap}{0.45\baselineskip}
\newlength{\exampleblockgap}
\setlength{\exampleblockgap}{0.65\baselineskip}
\newlength{\passageparagraphgap}
\setlength{\passageparagraphgap}{0.45\baselineskip}
\newenvironment{exampleblock}{\par\vspace{0.35\baselineskip}\begingroup\raggedright\setlength{\parskip}{0pt}\setlength{\parindent}{0pt}}{\par\endgroup\vspace{\exampleblockgap}}
\newcommand{\readinglineindent}[2]{%
\par\noindent%
\hangindent=\dimexpr\readinggutter+0.9em\relax%
\hangafter=1%
\makebox[\readinggutter][r]{\scriptsize\color{gray}#1}\hspace{0.9em}\hspace*{\readingfirstindent}#2%
\par}
\newcommand{\readingline}[2]{%
\par\noindent%
\hangindent=\dimexpr\readinggutter+0.9em\relax%
\hangafter=1%
\makebox[\readinggutter][r]{\scriptsize\color{gray}#1}\hspace{0.9em}#2%
\par}

\setlength{\parindent}{0pt}
\setlength{\parskip}{0.5em}
\setlist[enumerate]{topsep=0.25em,itemsep=0.18em,parsep=0pt,partopsep=0pt}
\setlength{\emergencystretch}{2em}

\begin{document}
\section*{TOEFL Practice 03 - Review}
Source of truth: DOCX only.
\section*{Listening Comprehension}
\subsection*{Part A}
DIRECTIONS: In Part A you will hear short conversations between two speakers. At the end of each conversation, a third speaker will ask a question about what the first two speakers said. Each conversation and each question will be spoken only one time. Therefore, you must listen carefully to understand what each speaker says. After you hear a conversation and the question, read the four choices and select the one that is the best answer to the question the speaker asked. Then, on your answer. sheet, find the number of the question and blacken the space that corresponds to the letter for the answer you have chosen. Blacken the space completely so that the letter inside the space does not show.\par
\begin{exampleblock}
Listen to the following example.\par
On the recording, you hear:\par
(Man) Does the car need to be filled?\par
(Woman) Mary stopped at the gas station on her way home.\par
(Narrator) What does the woman mean?\par
In your test book, you will read:\par
(A) Mary bought some food.\par
(B) Mary had car trouble.\par
(C) Mary went shopping.\par
(D) Mary bought some gas.\par
From the conversation you learn that Mary stopped at the gas station on her way home. The best answer to the question "Does the car need to be filled?" is (D), "Mary bought some gas." Therefore, the correct answer is (D).\par
Now let us begin Part A with question number 1.\par
\end{exampleblock}
\begin{exampleblock}
\textbf{Transkrip rekonstruksi AI (draft; bukan resmi)}\par
Man: How did Susan do on her driving test?\par
Woman: She failed it, and she was really surprised. She had been sure she would pass.\par
Narrator: What does the woman mean?\par
\end{exampleblock}
\noindent\textbf{1.}
\begin{enumerate}[label=(\Alph*),leftmargin=*,itemsep=0.2em]
\item She didn't expect to fail.
\item She went sailing.
\item She had her hearing examined.
\item She passed her driving test.
\end{enumerate}
\par
\begingroup
\setlength{\parskip}{0pt}
\setlength{\parindent}{0pt}
\textbf{Review}\par
Answer: A\par
Confidence: medium\par
Jawaban diambil dari kunci belajar yang diberikan pengguna dan belum diverifikasi terhadap audio/transkrip resmi.\par
\endgroup
\par
\vspace{0.25\baselineskip}
\begin{exampleblock}
\textbf{Transkrip rekonstruksi AI (draft; bukan resmi)}\par
Woman: Are we late for the movie?\par
Man: Not because of the drive. The real problem was finding a place to park; we drove around the lot for ages.\par
Narrator: What does the man mean?\par
\end{exampleblock}
\noindent\textbf{2.}
\begin{enumerate}[label=(\Alph*),leftmargin=*,itemsep=0.2em]
\item We have 40 minutes to spare.
\item We have to be there in 14 minutes.
\item We had a hard time parking the car.
\item The parking lot is around the corner.
\end{enumerate}
\par
\begingroup
\setlength{\parskip}{0pt}
\setlength{\parindent}{0pt}
\textbf{Review}\par
Answer: C\par
Confidence: medium\par
Jawaban diambil dari kunci belajar yang diberikan pengguna dan belum diverifikasi terhadap audio/transkrip resmi.\par
\endgroup
\par
\vspace{0.25\baselineskip}
\begin{exampleblock}
\textbf{Transkrip rekonstruksi AI (draft; bukan resmi)}\par
Woman: How are Tom's investments doing?\par
Man: Not very well. He says his current broker has made too many bad decisions, so he needs a new broker.\par
Narrator: What does the man say about Tom?\par
\end{exampleblock}
\noindent\textbf{3.}
\begin{enumerate}[label=(\Alph*),leftmargin=*,itemsep=0.2em]
\item I'm wearing his shoes.
\item He needs a new broker.
\item I broke my shoe strap.
\item The manager made investments.
\end{enumerate}
\par
\begingroup
\setlength{\parskip}{0pt}
\setlength{\parindent}{0pt}
\textbf{Review}\par
Answer: B\par
Confidence: medium\par
Jawaban diambil dari kunci belajar yang diberikan pengguna dan belum diverifikasi terhadap audio/transkrip resmi.\par
\endgroup
\par
\vspace{0.25\baselineskip}
\begin{exampleblock}
\textbf{Transkrip rekonstruksi AI (draft; bukan resmi)}\par
Woman: I'm here for Professor Harris. Is it all right if I go in?\par
Man: Sure. The door is open, and she is expecting you. You can go right in.\par
Narrator: What does the man tell the woman?\par
\end{exampleblock}
\noindent\textbf{4.}
\begin{enumerate}[label=(\Alph*),leftmargin=*,itemsep=0.2em]
\item She can go right in.
\item She can arrive by herself.
\item He can't let her in m y house.
\item He doesn't have a key for her.
\end{enumerate}
\par
\begingroup
\setlength{\parskip}{0pt}
\setlength{\parindent}{0pt}
\textbf{Review}\par
Answer: A\par
Confidence: medium\par
Jawaban diambil dari kunci belajar yang diberikan pengguna dan belum diverifikasi terhadap audio/transkrip resmi.\par
\endgroup
\par
\vspace{0.25\baselineskip}
\begin{exampleblock}
\textbf{Transkrip rekonstruksi AI (draft; bukan resmi)}\par
Man: Which ice cream flavor is the most popular these days?\par
Woman: Even with all the new flavors, vanilla is still the favorite.\par
Narrator: What does the woman mean?\par
\end{exampleblock}
\noindent\textbf{5.}
\begin{enumerate}[label=(\Alph*),leftmargin=*,itemsep=0.2em]
\item Seventy flavors are sold regularly.
\item Vanilla is the favorite flavor.
\item Ice cream was created in the past decade.
\item Vanilla accounts are 70 percent accurate.
\end{enumerate}
\par
\begingroup
\setlength{\parskip}{0pt}
\setlength{\parindent}{0pt}
\textbf{Review}\par
Answer: B\par
Confidence: medium\par
Jawaban diambil dari kunci belajar yang diberikan pengguna dan belum diverifikasi terhadap audio/transkrip resmi.\par
\endgroup
\par
\vspace{0.25\baselineskip}
\begin{exampleblock}
\textbf{Transkrip rekonstruksi AI (draft; bukan resmi)}\par
Woman: Could you give me a ride home? My car won't start again.\par
Man: Again? You really are having trouble with that car.\par
Narrator: What does the man mean?\par
\end{exampleblock}
\noindent\textbf{6.}
\begin{enumerate}[label=(\Alph*),leftmargin=*,itemsep=0.2em]
\item She should buy a new desk light.
\item She had to think of a new headline.
\item She is having trouble with her car.
\item She should talk to the new hires.
\end{enumerate}
\par
\begingroup
\setlength{\parskip}{0pt}
\setlength{\parindent}{0pt}
\textbf{Review}\par
Answer: C\par
Confidence: medium\par
Jawaban diambil dari kunci belajar yang diberikan pengguna dan belum diverifikasi terhadap audio/transkrip resmi.\par
\endgroup
\par
\vspace{0.25\baselineskip}
\begin{exampleblock}
\textbf{Transkrip rekonstruksi AI (draft; bukan resmi)}\par
Man: Did Mark get to see Joanna before she left?\par
Woman: No, he arrived just after she had gone. He missed her by only a few minutes.\par
Narrator: What does the woman mean?\par
\end{exampleblock}
\noindent\textbf{7.}
\begin{enumerate}[label=(\Alph*),leftmargin=*,itemsep=0.2em]
\item Mark missed Joanna.
\item Joe and Ann love canoeing.
\item Joanna left after 10 o'clock.
\item Mark called his brother.
\end{enumerate}
\par
\begingroup
\setlength{\parskip}{0pt}
\setlength{\parindent}{0pt}
\textbf{Review}\par
Answer: A\par
Confidence: medium\par
Jawaban diambil dari kunci belajar yang diberikan pengguna dan belum diverifikasi terhadap audio/transkrip resmi.\par
\endgroup
\par
\vspace{0.25\baselineskip}
\begin{exampleblock}
\textbf{Transkrip rekonstruksi AI (draft; bukan resmi)}\par
Woman: I heard the company changed its pay policy.\par
Man: Unfortunately, yes. Everyone's salary was reduced this month.\par
Narrator: What does the man mean?\par
\end{exampleblock}
\noindent\textbf{8.}
\begin{enumerate}[label=(\Alph*),leftmargin=*,itemsep=0.2em]
\item Their revenue has increased.
\item Their salaries were reduced.
\item The income has been cut off.
\item The bay has been divided.
\end{enumerate}
\par
\begingroup
\setlength{\parskip}{0pt}
\setlength{\parindent}{0pt}
\textbf{Review}\par
Answer: B\par
Confidence: medium\par
Jawaban diambil dari kunci belajar yang diberikan pengguna dan belum diverifikasi terhadap audio/transkrip resmi.\par
\endgroup
\par
\vspace{0.25\baselineskip}
\begin{exampleblock}
\textbf{Transkrip rekonstruksi AI (draft; bukan resmi)}\par
Man: How much did Mrs. Bailey pay for the flowers?\par
Woman: They cost fifty dollars, including the delivery charge.\par
Narrator: What does the woman say?\par
\end{exampleblock}
\noindent\textbf{9.}
\begin{enumerate}[label=(\Alph*),leftmargin=*,itemsep=0.2em]
\item Mrs. Bailey was cheated out of \$38.
\item Mrs. Bailey charged the flowers on her credit card.
\item Mrs. Bailey charges a lot for her flower delivery.
\item Mrs. Bailey paid \$50 for the flowers.
\end{enumerate}
\par
\begingroup
\setlength{\parskip}{0pt}
\setlength{\parindent}{0pt}
\textbf{Review}\par
Answer: D\par
Confidence: medium\par
Jawaban diambil dari kunci belajar yang diberikan pengguna dan belum diverifikasi terhadap audio/transkrip resmi.\par
\endgroup
\par
\vspace{0.25\baselineskip}
\begin{exampleblock}
\textbf{Transkrip rekonstruksi AI (draft; bukan resmi)}\par
Woman: Do you think the new proposal will be approved?\par
Man: It seems stronger than the one they rejected last month. I think this proposal has a better chance.\par
Narrator: What does the man imply?\par
\end{exampleblock}
\noindent\textbf{10.}
\begin{enumerate}[label=(\Alph*),leftmargin=*,itemsep=0.2em]
\item The new hospital works better than the old one.
\item The hospital was built six years ago but still looks good.
\item They built a hospital here six years ago.
\item This proposal has a better chance than the previous one.
\end{enumerate}
\par
\begingroup
\setlength{\parskip}{0pt}
\setlength{\parindent}{0pt}
\textbf{Review}\par
Answer: D\par
Confidence: medium\par
Jawaban diambil dari kunci belajar yang diberikan pengguna dan belum diverifikasi terhadap audio/transkrip resmi.\par
\endgroup
\par
\vspace{0.25\baselineskip}
\begin{exampleblock}
\textbf{Transkrip rekonstruksi AI (draft; bukan resmi)}\par
Man: Why is the bell ringing in the library?\par
Woman: It means the library is closing. Everyone using the library should leave now.\par
Narrator: What does the woman mean?\par
\end{exampleblock}
\noindent\textbf{11.}
\begin{enumerate}[label=(\Alph*),leftmargin=*,itemsep=0.2em]
\item The bell sounds like it's coming from the library.
\item The library has a sound strategy for patrons.
\item The library users should leave when the bell rings.
\item Parents shouldn't leave children in the library.
\end{enumerate}
\par
\begingroup
\setlength{\parskip}{0pt}
\setlength{\parindent}{0pt}
\textbf{Review}\par
Answer: C\par
Confidence: medium\par
Jawaban diambil dari kunci belajar yang diberikan pengguna dan belum diverifikasi terhadap audio/transkrip resmi.\par
\endgroup
\par
\vspace{0.25\baselineskip}
\begin{exampleblock}
\textbf{Transkrip rekonstruksi AI (draft; bukan resmi)}\par
Woman: Should we meet in front of the music hall, or should we go together?\par
Man: We don't need to go together. We can go to the music hall separately and meet inside.\par
Narrator: What does the man suggest?\par
\end{exampleblock}
\noindent\textbf{12.}
\begin{enumerate}[label=(\Alph*),leftmargin=*,itemsep=0.2em]
\item We can go to the music hall separately.
\item We don't have to stand during the concert.
\item Let's have a meat dish after the concert.
\item We can have the meeting in the hall.
\end{enumerate}
\par
\begingroup
\setlength{\parskip}{0pt}
\setlength{\parindent}{0pt}
\textbf{Review}\par
Answer: A\par
Confidence: medium\par
Jawaban diambil dari kunci belajar yang diberikan pengguna dan belum diverifikasi terhadap audio/transkrip resmi.\par
\endgroup
\par
\vspace{0.25\baselineskip}
\begin{exampleblock}
\textbf{Transkrip rekonstruksi AI (draft; bukan resmi)}\par
Man: Was Clara ready for her presentation?\par
Woman: Definitely. She spent hours preparing the slides and checking every detail.\par
Narrator: What does the woman mean?\par
\end{exampleblock}
\noindent\textbf{13.}
\begin{enumerate}[label=(\Alph*),leftmargin=*,itemsep=0.2em]
\item She prepared the slides carefully.
\item She was sure she did well at the presentation.
\item She presented the slides to the sales personnel.
\item She checked the slides out of the sales office.
\end{enumerate}
\par
\begingroup
\setlength{\parskip}{0pt}
\setlength{\parindent}{0pt}
\textbf{Review}\par
Answer: A\par
Confidence: medium\par
Jawaban diambil dari kunci belajar yang diberikan pengguna dan belum diverifikasi terhadap audio/transkrip resmi.\par
\endgroup
\par
\vspace{0.25\baselineskip}
\begin{exampleblock}
\textbf{Transkrip rekonstruksi AI (draft; bukan resmi)}\par
Woman: How many friends did you invite for the break?\par
Man: I invited eight friends, but I'm not sure all of them can come.\par
Narrator: What does the man say?\par
\end{exampleblock}
\noindent\textbf{14.}
\begin{enumerate}[label=(\Alph*),leftmargin=*,itemsep=0.2em]
\item Four people are coming.
\item Will you come for the break?
\item Would you like to visit us?
\item Eight friends are invited.
\end{enumerate}
\par
\begingroup
\setlength{\parskip}{0pt}
\setlength{\parindent}{0pt}
\textbf{Review}\par
Answer: D\par
Confidence: medium\par
Jawaban diambil dari kunci belajar yang diberikan pengguna dan belum diverifikasi terhadap audio/transkrip resmi.\par
\endgroup
\par
\vspace{0.25\baselineskip}
\begin{exampleblock}
\textbf{Transkrip rekonstruksi AI (draft; bukan resmi)}\par
Man: What is Margaret planning to do now that her children are older?\par
Woman: She is thinking about going back to school and taking some college classes.\par
Narrator: What does the woman say about Margaret?\par
\end{exampleblock}
\noindent\textbf{15.}
\begin{enumerate}[label=(\Alph*),leftmargin=*,itemsep=0.2em]
\item Margaret already has a college degree.
\item Margaret is considering taking classes.
\item Margaret's children are away at college.
\item Margaret's children are outside.
\end{enumerate}
\par
\begingroup
\setlength{\parskip}{0pt}
\setlength{\parindent}{0pt}
\textbf{Review}\par
Answer: B\par
Confidence: medium\par
Jawaban diambil dari kunci belajar yang diberikan pengguna dan belum diverifikasi terhadap audio/transkrip resmi.\par
\endgroup
\par
\vspace{0.25\baselineskip}
\begin{exampleblock}
\textbf{Transkrip rekonstruksi AI (draft; bukan resmi)}\par
Woman: Ralph looked worried when the boys left. What do you think he was thinking?\par
Man: It's hard to tell with Ralph. He never says much about what is on his mind.\par
Narrator: What does the man imply?\par
\end{exampleblock}
\noindent\textbf{16.}
\begin{enumerate}[label=(\Alph*),leftmargin=*,itemsep=0.2em]
\item It's difficult to say what Ralph is thinking.
\item The boys are doing well, and Ralph is not worried
\item Ralph only seems worried, but he is really not
\item Ralph is always thinking about the boys.
\end{enumerate}
\par
\begingroup
\setlength{\parskip}{0pt}
\setlength{\parindent}{0pt}
\textbf{Review}\par
Answer: A\par
Confidence: medium\par
Jawaban diambil dari kunci belajar yang diberikan pengguna dan belum diverifikasi terhadap audio/transkrip resmi.\par
\endgroup
\par
\vspace{0.25\baselineskip}
\begin{exampleblock}
\textbf{Transkrip rekonstruksi AI (draft; bukan resmi)}\par
Man: Good morning. I'm calling about my car. Is it ready yet?\par
Woman: Let me check with the mechanic. May I have your name and the repair order number?\par
Narrator: Who is the woman most likely to be?\par
\end{exampleblock}
\noindent\textbf{17.}
\begin{enumerate}[label=(\Alph*),leftmargin=*,itemsep=0.2em]
\item A receptionist in a car repair shop.
\item An operator at a telephone company.
\item A veterinarian at an animal clinic.
\item A secretary in an insurance company.
\end{enumerate}
\par
\begingroup
\setlength{\parskip}{0pt}
\setlength{\parindent}{0pt}
\textbf{Review}\par
Answer: A\par
Confidence: medium\par
Jawaban diambil dari kunci belajar yang diberikan pengguna dan belum diverifikasi terhadap audio/transkrip resmi.\par
\endgroup
\par
\vspace{0.25\baselineskip}
\begin{exampleblock}
\textbf{Transkrip rekonstruksi AI (draft; bukan resmi)}\par
Woman: Is the grocery store still open this late?\par
Man: Yes. It stays open until midnight, so we can still go shopping now.\par
Narrator: What does the man mean?\par
\end{exampleblock}
\noindent\textbf{18.}
\begin{enumerate}[label=(\Alph*),leftmargin=*,itemsep=0.2em]
\item The grocery store has an odd schedule.
\item Arthur stores things at his place.
\item It's possible to go shopping at that hour.
\item Only strange people go to the store so late.
\end{enumerate}
\par
\begingroup
\setlength{\parskip}{0pt}
\setlength{\parindent}{0pt}
\textbf{Review}\par
Answer: C\par
Confidence: medium\par
Jawaban diambil dari kunci belajar yang diberikan pengguna dan belum diverifikasi terhadap audio/transkrip resmi.\par
\endgroup
\par
\vspace{0.25\baselineskip}
\begin{exampleblock}
\textbf{Transkrip rekonstruksi AI (draft; bukan resmi)}\par
Woman: Why did Paul stop at the service station?\par
Man: He had almost no gas left. In fact, he was afraid he might not make it there.\par
Narrator: What does the man imply about Paul?\par
\end{exampleblock}
\noindent\textbf{19.}
\begin{enumerate}[label=(\Alph*),leftmargin=*,itemsep=0.2em]
\item He ran out of the station.
\item He didn't have enough gas.
\item He has been thinking a lot.
\item He pulled over to buy gas.
\end{enumerate}
\par
\begingroup
\setlength{\parskip}{0pt}
\setlength{\parindent}{0pt}
\textbf{Review}\par
Answer: B\par
Confidence: medium\par
Jawaban diambil dari kunci belajar yang diberikan pengguna dan belum diverifikasi terhadap audio/transkrip resmi.\par
\endgroup
\par
\vspace{0.25\baselineskip}
\begin{exampleblock}
\textbf{Transkrip rekonstruksi AI (draft; bukan resmi)}\par
Man: I'm confused about stress in English words. Can a word have two main stresses?\par
Woman: Usually not. Most words have only one primary stress.\par
Narrator: What does the woman mean?\par
\end{exampleblock}
\noindent\textbf{20.}
\begin{enumerate}[label=(\Alph*),leftmargin=*,itemsep=0.2em]
\item Two-syllable words are common.
\item Stress can affect pronunciation.
\item Two-part words are very rare.
\item Most words have only one stress.
\end{enumerate}
\par
\begingroup
\setlength{\parskip}{0pt}
\setlength{\parindent}{0pt}
\textbf{Review}\par
Answer: D\par
Confidence: medium\par
Jawaban diambil dari kunci belajar yang diberikan pengguna dan belum diverifikasi terhadap audio/transkrip resmi.\par
\endgroup
\par
\vspace{0.25\baselineskip}
\begin{exampleblock}
\textbf{Transkrip rekonstruksi AI (draft; bukan resmi)}\par
Woman: Where should I put these dishes?\par
Man: Put them in the cupboard above the sink, next to the plates.\par
Narrator: Where are the speakers?\par
\end{exampleblock}
\noindent\textbf{21.}
\begin{enumerate}[label=(\Alph*),leftmargin=*,itemsep=0.2em]
\item In a kitchen
\item In a living room
\item In a bathroom
\item In a hallway
\end{enumerate}
\par
\begingroup
\setlength{\parskip}{0pt}
\setlength{\parindent}{0pt}
\textbf{Review}\par
Answer: A\par
Confidence: medium\par
Jawaban diambil dari kunci belajar yang diberikan pengguna dan belum diverifikasi terhadap audio/transkrip resmi.\par
\endgroup
\par
\vspace{0.25\baselineskip}
\begin{exampleblock}
\textbf{Transkrip rekonstruksi AI (draft; bukan resmi)}\par
Man: Where will Helen stay while her apartment is being painted?\par
Woman: She said she would stay at her sister's place for a few days.\par
Narrator: Where will Helen stay?\par
\end{exampleblock}
\noindent\textbf{22.}
\begin{enumerate}[label=(\Alph*),leftmargin=*,itemsep=0.2em]
\item At her father's
\item At her sister's
\item At her own home
\item At her children's
\end{enumerate}
\par
\begingroup
\setlength{\parskip}{0pt}
\setlength{\parindent}{0pt}
\textbf{Review}\par
Answer: B\par
Confidence: medium\par
Jawaban diambil dari kunci belajar yang diberikan pengguna dan belum diverifikasi terhadap audio/transkrip resmi.\par
\endgroup
\par
\vspace{0.25\baselineskip}
\begin{exampleblock}
\textbf{Transkrip rekonstruksi AI (draft; bukan resmi)}\par
Woman: I need something for this cough.\par
Man: The pharmacist can help you choose a suitable medicine. The counter is over there.\par
Narrator: Where are the speakers?\par
\end{exampleblock}
\noindent\textbf{23.}
\begin{enumerate}[label=(\Alph*),leftmargin=*,itemsep=0.2em]
\item At a post office
\item At a bank
\item At a pharmacy
\item At a deli
\end{enumerate}
\par
\begingroup
\setlength{\parskip}{0pt}
\setlength{\parindent}{0pt}
\textbf{Review}\par
Answer: C\par
Confidence: medium\par
Jawaban diambil dari kunci belajar yang diberikan pengguna dan belum diverifikasi terhadap audio/transkrip resmi.\par
\endgroup
\par
\vspace{0.25\baselineskip}
\begin{exampleblock}
\textbf{Transkrip rekonstruksi AI (draft; bukan resmi)}\par
Man: The mechanic said your car should be ready by three-thirty.\par
Woman: Three-thirty? I doubt it. They told me earlier that the work would take until at least five.\par
Narrator: What does the woman mean?\par
\end{exampleblock}
\noindent\textbf{24.}
\begin{enumerate}[label=(\Alph*),leftmargin=*,itemsep=0.2em]
\item She thinks her car will be ready before 5:00.
\item She is surprised that her car is not finished.
\item She doesn't believe her car will be ready at 3:30.
\item She wants the man to give her a ride to the shop.
\end{enumerate}
\par
\begingroup
\setlength{\parskip}{0pt}
\setlength{\parindent}{0pt}
\textbf{Review}\par
Answer: C\par
Confidence: medium\par
Jawaban diambil dari kunci belajar yang diberikan pengguna dan belum diverifikasi terhadap audio/transkrip resmi.\par
\endgroup
\par
\vspace{0.25\baselineskip}
\begin{exampleblock}
\textbf{Transkrip rekonstruksi AI (draft; bukan resmi)}\par
Woman: Why does John look so uncomfortable?\par
Man: He had two large bowls of soup and then ate a full dinner. I think he ate too much.\par
Narrator: What does the man imply about John?\par
\end{exampleblock}
\noindent\textbf{25.}
\begin{enumerate}[label=(\Alph*),leftmargin=*,itemsep=0.2em]
\item He is tired.
\item He ate too much.
\item He is out of breath.
\item He made soup.
\end{enumerate}
\par
\begingroup
\setlength{\parskip}{0pt}
\setlength{\parindent}{0pt}
\textbf{Review}\par
Answer: B\par
Confidence: medium\par
Jawaban diambil dari kunci belajar yang diberikan pengguna dan belum diverifikasi terhadap audio/transkrip resmi.\par
\endgroup
\par
\vspace{0.25\baselineskip}
\begin{exampleblock}
\textbf{Transkrip rekonstruksi AI (draft; bukan resmi)}\par
Man: Did Carol mind that the Nixons came with us?\par
Woman: No, not at all. She said it was fine.\par
Narrator: What does the woman mean?\par
\end{exampleblock}
\noindent\textbf{26.}
\begin{enumerate}[label=(\Alph*),leftmargin=*,itemsep=0.2em]
\item She couldn't be bothered.
\item She invited the Nixons.
\item She didn't mind.
\item She has a brother.
\end{enumerate}
\par
\begingroup
\setlength{\parskip}{0pt}
\setlength{\parindent}{0pt}
\textbf{Review}\par
Answer: C\par
Confidence: medium\par
Jawaban diambil dari kunci belajar yang diberikan pengguna dan belum diverifikasi terhadap audio/transkrip resmi.\par
\endgroup
\par
\vspace{0.25\baselineskip}
\begin{exampleblock}
\textbf{Transkrip rekonstruksi AI (draft; bukan resmi)}\par
Man: Did you enjoy your evening with Nick?\par
Woman: Not really. All he talked about was tennis, and after a while I got bored.\par
Narrator: What does the woman imply?\par
\end{exampleblock}
\noindent\textbf{27.}
\begin{enumerate}[label=(\Alph*),leftmargin=*,itemsep=0.2em]
\item Nick plays tennis all the time.
\item The woman went out with Nick.
\item The town has many nightclubs.
\item The woman finds Nick boring.
\end{enumerate}
\par
\begingroup
\setlength{\parskip}{0pt}
\setlength{\parindent}{0pt}
\textbf{Review}\par
Answer: D\par
Confidence: medium\par
Jawaban diambil dari kunci belajar yang diberikan pengguna dan belum diverifikasi terhadap audio/transkrip resmi.\par
\endgroup
\par
\vspace{0.25\baselineskip}
\begin{exampleblock}
\textbf{Transkrip rekonstruksi AI (draft; bukan resmi)}\par
Woman: Is the introductory economics class very large?\par
Man: Yes, it is. Two hundred students are enrolled in it this semester.\par
Narrator: What does the man say about the class?\par
\end{exampleblock}
\noindent\textbf{28.}
\begin{enumerate}[label=(\Alph*),leftmargin=*,itemsep=0.2em]
\item Four hundred students will take it.
\item Two hundred students are in it.
\item It's restricted to the waiting list.
\item Its instructor has strict rules.
\end{enumerate}
\par
\begingroup
\setlength{\parskip}{0pt}
\setlength{\parindent}{0pt}
\textbf{Review}\par
Answer: B\par
Confidence: medium\par
Jawaban diambil dari kunci belajar yang diberikan pengguna dan belum diverifikasi terhadap audio/transkrip resmi.\par
\endgroup
\par
\vspace{0.25\baselineskip}
\begin{exampleblock}
\textbf{Transkrip rekonstruksi AI (draft; bukan resmi)}\par
Man: Can you shorten these pants by Friday?\par
Woman: Yes. I'll measure the length now and have the alterations finished by then.\par
Narrator: Where does this conversation most likely take place?\par
\end{exampleblock}
\noindent\textbf{29.}
\begin{enumerate}[label=(\Alph*),leftmargin=*,itemsep=0.2em]
\item At a laundromat
\item At a printmaker's
\item In an appliance store
\item In a tailor's shop
\end{enumerate}
\par
\begingroup
\setlength{\parskip}{0pt}
\setlength{\parindent}{0pt}
\textbf{Review}\par
Answer: D\par
Confidence: medium\par
Jawaban diambil dari kunci belajar yang diberikan pengguna dan belum diverifikasi terhadap audio/transkrip resmi.\par
\endgroup
\par
\vspace{0.25\baselineskip}
\begin{exampleblock}
\textbf{Transkrip rekonstruksi AI (draft; bukan resmi)}\par
Woman: Please fasten your seat belt and return your tray table to its upright position. We'll be landing shortly.\par
Man: Of course. Thank you.\par
Narrator: Who is the woman most likely to be?\par
\end{exampleblock}
\noindent\textbf{30.}
\begin{enumerate}[label=(\Alph*),leftmargin=*,itemsep=0.2em]
\item A doctor
\item A flight attendant
\item A hospital nurse
\item A pilot
\end{enumerate}
\par
\begingroup
\setlength{\parskip}{0pt}
\setlength{\parindent}{0pt}
\textbf{Review}\par
Answer: B\par
Confidence: medium\par
Jawaban diambil dari kunci belajar yang diberikan pengguna dan belum diverifikasi terhadap audio/transkrip resmi.\par
\endgroup
\par
\vspace{0.25\baselineskip}
\subsection*{Part B}
DIRECTIONS: In this part of the test, you will hear longer conversations. After each conversation, you will hear several questions. The conversations and questions will not be repeated. After you hear a question, read the four possible answers in. your test book and select. The best answer. Then, on your answer sheet, find the number of the question and fill in the space that corresponds to the letter of the answer you have chosen. Remember, you are not allowed to take notes or write in your test book.\par
\begin{exampleblock}
Listen to the following example:\par
You will hear:\par
You will read:\par
(A) He has changed jobs.\par
(B) He has two children.\par
(C) He has two jobs.\par
(D) He is looking for a job.\par
From the conversation you learn that Tom has taken an additional job. The best answer to the question "Why is Tom tired?" is (C), "He has two jobs." Therefore the correct answer is (C).\par
\end{exampleblock}
\begin{exampleblock}
\textbf{Transkrip rekonstruksi AI (draft; bukan resmi)}\par
Guide: Welcome to the marine mammal exhibit. Today we're going to talk about dolphins and some of their special features.\par
Student: I know dolphins live in water, but they aren't fish, right?\par
Guide: That's right. Dolphins are mammals. They breathe air, give birth to live young, and nurse their babies. They are also very social and intelligent animals.\par
Student: How long do baby dolphins stay with their mothers?\par
Guide: Usually for almost a year. During that time, the mother protects the calf and teaches it how to swim, feed, and stay with the group.\par
Student: Do dolphins communicate with each other?\par
Guide: Yes. They use a variety of sounds, such as whistles and clicks. These sounds help them identify one another, stay together, and sometimes warn the group of danger.\par
Student: That's amazing. I didn't realize they relied so much on sound.\par
Guide: Sound is extremely important to dolphins, especially because they often need to communicate underwater where visual contact may be limited.\par
\end{exampleblock}
\noindent\textbf{31.}
\begin{enumerate}[label=(\Alph*),leftmargin=*,itemsep=0.2em]
\item The life of sociable creatures
\item Special features of dolphins
\item The differences between fish and dolphins
\item Communication between groups of dolphins
\end{enumerate}
\par
\begingroup
\setlength{\parskip}{0pt}
\setlength{\parindent}{0pt}
\textbf{Review}\par
Answer: B\par
Confidence: medium\par
Jawaban diambil dari kunci belajar yang diberikan pengguna dan belum diverifikasi terhadap audio/transkrip resmi.\par
\endgroup
\par
\vspace{0.25\baselineskip}
\noindent\textbf{32.}
\begin{enumerate}[label=(\Alph*),leftmargin=*,itemsep=0.2em]
\item In a museum
\item On a seashore
\item In a laboratory
\item At a zoo
\end{enumerate}
\par
\begingroup
\setlength{\parskip}{0pt}
\setlength{\parindent}{0pt}
\textbf{Review}\par
Answer: D\par
Confidence: medium\par
Jawaban diambil dari kunci belajar yang diberikan pengguna dan belum diverifikasi terhadap audio/transkrip resmi.\par
\endgroup
\par
\vspace{0.25\baselineskip}
\noindent\textbf{33.}
\begin{enumerate}[label=(\Alph*),leftmargin=*,itemsep=0.2em]
\item For about 2 months
\item For almost 12 months
\item Approximately 2 years
\item Approximately 20 years
\end{enumerate}
\par
\begingroup
\setlength{\parskip}{0pt}
\setlength{\parindent}{0pt}
\textbf{Review}\par
Answer: B\par
Confidence: medium\par
Jawaban diambil dari kunci belajar yang diberikan pengguna dan belum diverifikasi terhadap audio/transkrip resmi.\par
\endgroup
\par
\vspace{0.25\baselineskip}
\noindent\textbf{34.}
\begin{enumerate}[label=(\Alph*),leftmargin=*,itemsep=0.2em]
\item By visual contact
\item By sounds
\item Through warning one another
\item With the help of scientists
\end{enumerate}
\par
\begingroup
\setlength{\parskip}{0pt}
\setlength{\parindent}{0pt}
\textbf{Review}\par
Answer: B\par
Confidence: medium\par
Jawaban diambil dari kunci belajar yang diberikan pengguna dan belum diverifikasi terhadap audio/transkrip resmi.\par
\endgroup
\par
\vspace{0.25\baselineskip}
\begin{exampleblock}
\textbf{Transkrip rekonstruksi AI (draft; bukan resmi)}\par
Student: Professor, I use a dictionary all the time, but sometimes it seems to make reading more confusing instead of easier.\par
Professor: That can happen. Dictionaries are very useful, but students need to understand both their benefits and their pitfalls.\par
Student: What do you mean by pitfalls?\par
Professor: A dictionary may give several definitions for one word, and not every definition fits the sentence you are reading. If you do not understand the word's usage in context, you may choose the wrong meaning.\par
Student: So dictionaries do not solve everything automatically.\par
Professor: Exactly. There are also different kinds of dictionaries. General dictionaries are useful for everyday vocabulary, while specialized dictionaries are designed for fields like medicine, law, engineering, or linguistics.\par
Student: Did you learn that from studying vocabulary?\par
Professor: I learned it most clearly when I taught a reading class. Students often had dictionaries, but they still misunderstood passages because they looked up words without considering grammar, usage, or context.\par
Student: So a dictionary helps only if the learner already knows how to use it effectively.\par
Professor: That's the point. A dictionary is a tool, not a substitute for understanding word usage.\par
\end{exampleblock}
\noindent\textbf{35.}
\begin{enumerate}[label=(\Alph*),leftmargin=*,itemsep=0.2em]
\item The contents and structure of dictionaries
\item The benefits and pitfalls of dictionary use
\item The acquisition of words in dictionaries
\item The various applications of dictionary entries
\end{enumerate}
\par
\begingroup
\setlength{\parskip}{0pt}
\setlength{\parindent}{0pt}
\textbf{Review}\par
Answer: B\par
Confidence: medium\par
Jawaban diambil dari kunci belajar yang diberikan pengguna dan belum diverifikasi terhadap audio/transkrip resmi.\par
\endgroup
\par
\vspace{0.25\baselineskip}
\noindent\textbf{36.}
\begin{enumerate}[label=(\Alph*),leftmargin=*,itemsep=0.2em]
\item Formal and technical
\item General and specialized
\item Small and unabridged
\item American and English
\end{enumerate}
\par
\begingroup
\setlength{\parskip}{0pt}
\setlength{\parindent}{0pt}
\textbf{Review}\par
Answer: B\par
Confidence: medium\par
Jawaban diambil dari kunci belajar yang diberikan pengguna dan belum diverifikasi terhadap audio/transkrip resmi.\par
\endgroup
\par
\vspace{0.25\baselineskip}
\noindent\textbf{37.}
\begin{enumerate}[label=(\Alph*),leftmargin=*,itemsep=0.2em]
\item Took a writing class
\item Used a technical dictionary
\item Studied many new words
\item Taught a reading class
\end{enumerate}
\par
\begingroup
\setlength{\parskip}{0pt}
\setlength{\parindent}{0pt}
\textbf{Review}\par
Answer: D\par
Confidence: medium\par
Jawaban diambil dari kunci belajar yang diberikan pengguna dan belum diverifikasi terhadap audio/transkrip resmi.\par
\endgroup
\par
\vspace{0.25\baselineskip}
\noindent\textbf{38.}
\begin{enumerate}[label=(\Alph*),leftmargin=*,itemsep=0.2em]
\item Learners need to abandon dictionaries when their vocabulary increases.
\item Teachers think that students shouldn't always depend on dictionaries.
\item One can't use a dictionary effectively without knowing word usage.
\item Although dictionaries contain many entries, some word definitions can be irrelevant.
\end{enumerate}
\par
\begingroup
\setlength{\parskip}{0pt}
\setlength{\parindent}{0pt}
\textbf{Review}\par
Answer: C\par
Confidence: medium\par
Jawaban diambil dari kunci belajar yang diberikan pengguna dan belum diverifikasi terhadap audio/transkrip resmi.\par
\endgroup
\par
\vspace{0.25\baselineskip}
\subsection*{Part C}
DIRECTIONS: In Part C you will hear short lectures and extended conversations. At the end of each, you will be asked several questions. Each lecture or conversation and each question will be spoken only one time. For this reason, you must listen carefully to understand what each speaker says. After you hear a question, read the four choices and select the one that best answers the question the speaker asked. Then, on your answer sheet, find the number of the question and blacken the space that contains the letter for the answer you have chosen. Answer all questions according to what is stated or implied in the lecture or conversation.\par
\begin{exampleblock}
Listen to this sample talk.\par
You will hear:\par
Now listen, to the following example.\par
You will hear:\par
You will read:\par
(A) By cars and carriages\par
(B) By bicycles, trains, and carriages\par
(C) On foot and by boat\par
(D) On board ships and trains\par
The best answer to the question "According to the speaker, how did people travel before the invention of the automobile?" is (B), "By bicycles, trains, and carriages." Therefore, the correct answer is (B). You are not allowed to make notes during the test.\par
\end{exampleblock}
\begin{exampleblock}
\textbf{Transkrip rekonstruksi AI (draft; bukan resmi)}\par
Speaker: Today we will discuss Lake Ontario, one of the five Great Lakes of North America. Although Lake Ontario is important historically and geographically, it is the smallest of the Great Lakes. It is located east of Lake Erie and is connected to the Atlantic Ocean through the St. Lawrence River system.\par
Lake Ontario is smaller and generally has less shipping traffic than some of the other Great Lakes, especially compared with those more heavily used for large-scale industrial transportation. Its length is about 193 miles, and its shoreline is approximately 480 miles.\par
The lake also has an important effect on the surrounding region's climate. Because large bodies of water heat and cool more slowly than land, Lake Ontario helps moderate weather patterns near its shores. This influence can reduce temperature extremes and affect local seasonal conditions.\par
\end{exampleblock}
\noindent\textbf{39.}
\begin{enumerate}[label=(\Alph*),leftmargin=*,itemsep=0.2em]
\item It's the northern.
\item It's the southern.
\item It's the deepest.
\item It's the smallest.
\end{enumerate}
\par
\begingroup
\setlength{\parskip}{0pt}
\setlength{\parindent}{0pt}
\textbf{Review}\par
Answer: D\par
Confidence: medium\par
Jawaban diambil dari kunci belajar yang diberikan pengguna dan belum diverifikasi terhadap audio/transkrip resmi.\par
\endgroup
\par
\vspace{0.25\baselineskip}
\noindent\textbf{40.}
\begin{enumerate}[label=(\Alph*),leftmargin=*,itemsep=0.2em]
\item It's heavily traveled all year round
\item It's not navigable for large boats.
\item It has less traffic than the other Great Lakes.
\item Its shores are inaccessible for landing.
\end{enumerate}
\par
\begingroup
\setlength{\parskip}{0pt}
\setlength{\parindent}{0pt}
\textbf{Review}\par
Answer: C\par
Confidence: medium\par
Jawaban diambil dari kunci belajar yang diberikan pengguna dan belum diverifikasi terhadap audio/transkrip resmi.\par
\endgroup
\par
\vspace{0.25\baselineskip}
\noindent\textbf{41.}
\begin{enumerate}[label=(\Alph*),leftmargin=*,itemsep=0.2em]
\item 53 miles
\item 193 miles
\item 480 miles
\item 7,500 miles
\end{enumerate}
\par
\begingroup
\setlength{\parskip}{0pt}
\setlength{\parindent}{0pt}
\textbf{Review}\par
Answer: C\par
Confidence: medium\par
Jawaban diambil dari kunci belajar yang diberikan pengguna dan belum diverifikasi terhadap audio/transkrip resmi.\par
\endgroup
\par
\vspace{0.25\baselineskip}
\noindent\textbf{42.}
\begin{enumerate}[label=(\Alph*),leftmargin=*,itemsep=0.2em]
\item Its water is used for irrigation of trees.
\item It moderates the weather patterns in the region.
\item It cools the air in the surrounding areas.
\item It has a considerable capacity to store water.
\end{enumerate}
\par
\begingroup
\setlength{\parskip}{0pt}
\setlength{\parindent}{0pt}
\textbf{Review}\par
Answer: B\par
Confidence: medium\par
Jawaban diambil dari kunci belajar yang diberikan pengguna dan belum diverifikasi terhadap audio/transkrip resmi.\par
\endgroup
\par
\vspace{0.25\baselineskip}
\begin{exampleblock}
\textbf{Transkrip rekonstruksi AI (draft; bukan resmi)}\par
Speaker: Edgar Allan Poe, one of the most important figures in American literature, was born in Boston in 1809. After the death of his mother and the disappearance of his father, Poe was taken in by John and Frances Allan of Richmond, Virginia. The Allans were a well-to-do family, and Poe grew up with access to education and social opportunities that many children did not have.\par
When he was a child, Poe traveled with the Allan family to Britain, where he attended school for several years. He returned to the United States at about the age of eleven and later continued his studies in Richmond. Poe was considered a competent student, especially in languages and literature.\par
As a young man, he entered the University of Virginia. However, his time there was troubled. Poe became involved in gambling and accumulated debts. This caused serious conflict with John Allan and eventually contributed to Poe leaving the university.\par
\end{exampleblock}
\noindent\textbf{43.}
\begin{enumerate}[label=(\Alph*),leftmargin=*,itemsep=0.2em]
\item England
\item Scotland
\item Richmond
\item Boston
\end{enumerate}
\par
\begingroup
\setlength{\parskip}{0pt}
\setlength{\parindent}{0pt}
\textbf{Review}\par
Answer: D\par
Confidence: medium\par
Jawaban diambil dari kunci belajar yang diberikan pengguna dan belum diverifikasi terhadap audio/transkrip resmi.\par
\endgroup
\par
\vspace{0.25\baselineskip}
\noindent\textbf{44.}
\begin{enumerate}[label=(\Alph*),leftmargin=*,itemsep=0.2em]
\item They were well-to-do.
\item They didn't love young Poe.
\item They were English.
\item They abandoned him.
\end{enumerate}
\par
\begingroup
\setlength{\parskip}{0pt}
\setlength{\parindent}{0pt}
\textbf{Review}\par
Answer: A\par
Confidence: medium\par
Jawaban diambil dari kunci belajar yang diberikan pengguna dan belum diverifikasi terhadap audio/transkrip resmi.\par
\endgroup
\par
\vspace{0.25\baselineskip}
\noindent\textbf{45.}
\begin{enumerate}[label=(\Alph*),leftmargin=*,itemsep=0.2em]
\item 9
\item 11
\item 15
\item 18
\end{enumerate}
\par
\begingroup
\setlength{\parskip}{0pt}
\setlength{\parindent}{0pt}
\textbf{Review}\par
Answer: B\par
Confidence: medium\par
Jawaban diambil dari kunci belajar yang diberikan pengguna dan belum diverifikasi terhadap audio/transkrip resmi.\par
\endgroup
\par
\vspace{0.25\baselineskip}
\noindent\textbf{46.}
\begin{enumerate}[label=(\Alph*),leftmargin=*,itemsep=0.2em]
\item He was a competent student.
\item He preferred studying in the academy.
\item He did not do well in his academics.
\item He decided to travel to Scotland instead.
\end{enumerate}
\par
\begingroup
\setlength{\parskip}{0pt}
\setlength{\parindent}{0pt}
\textbf{Review}\par
Answer: A\par
Confidence: medium\par
Jawaban diambil dari kunci belajar yang diberikan pengguna dan belum diverifikasi terhadap audio/transkrip resmi.\par
\endgroup
\par
\vspace{0.25\baselineskip}
\noindent\textbf{47.}
\begin{enumerate}[label=(\Alph*),leftmargin=*,itemsep=0.2em]
\item John Allan could not afford Poe's education.
\item Poe became involved in gambling.
\item John Allan had to leave for England.
\item Poe decided to become a writer.
\end{enumerate}
\par
\begingroup
\setlength{\parskip}{0pt}
\setlength{\parindent}{0pt}
\textbf{Review}\par
Answer: B\par
Confidence: medium\par
Jawaban diambil dari kunci belajar yang diberikan pengguna dan belum diverifikasi terhadap audio/transkrip resmi.\par
\endgroup
\par
\vspace{0.25\baselineskip}
\begin{exampleblock}
\textbf{Transkrip rekonstruksi AI (draft; bukan resmi)}\par
Speaker: In marketing, it is useful to distinguish advertising from publicity. Advertising consists of presentations or messages designed to promote a product, service, organization, or idea. These messages are usually sent through mass media so that they can reach many people at the same time.\par
One key characteristic of advertising is that it is paid for. A company or sponsor buys space or time in a medium, such as television, radio, magazines, websites, or billboards. Because the sponsor pays for the message, the sponsor can usually control the content and timing of the presentation.\par
Publicity is different. Publicity may also bring attention to a product or organization, but it is not usually purchased in the same direct way. For example, a news story about a company may create publicity, but the company does not necessarily pay for the story as it would pay for an advertisement.\par
\end{exampleblock}
\noindent\textbf{48.}
\begin{enumerate}[label=(\Alph*),leftmargin=*,itemsep=0.2em]
\item It is nonpersonal.
\item It consists of presentations.
\item It cannot be bought.
\item It includes hamburgers and soft drinks.
\end{enumerate}
\par
\begingroup
\setlength{\parskip}{0pt}
\setlength{\parindent}{0pt}
\textbf{Review}\par
Answer: B\par
Confidence: medium\par
Jawaban diambil dari kunci belajar yang diberikan pengguna dan belum diverifikasi terhadap audio/transkrip resmi.\par
\endgroup
\par
\vspace{0.25\baselineskip}
\noindent\textbf{49.}
\begin{enumerate}[label=(\Alph*),leftmargin=*,itemsep=0.2em]
\item Product sponsors
\item Many people
\item TV viewers
\item Buyers of media time
\end{enumerate}
\par
\begingroup
\setlength{\parskip}{0pt}
\setlength{\parindent}{0pt}
\textbf{Review}\par
Answer: B\par
Confidence: medium\par
Jawaban diambil dari kunci belajar yang diberikan pengguna dan belum diverifikasi terhadap audio/transkrip resmi.\par
\endgroup
\par
\vspace{0.25\baselineskip}
\noindent\textbf{50.}
\begin{enumerate}[label=(\Alph*),leftmargin=*,itemsep=0.2em]
\item Advertising is paid for.
\item Advertising can be subject to skepticism.
\item Publicity reaches a greater audience.
\item Publicity usually benefits political leaders.
\end{enumerate}
\par
\begingroup
\setlength{\parskip}{0pt}
\setlength{\parindent}{0pt}
\textbf{Review}\par
Answer: A\par
Confidence: medium\par
Jawaban diambil dari kunci belajar yang diberikan pengguna dan belum diverifikasi terhadap audio/transkrip resmi.\par
\endgroup
\par
\vspace{0.25\baselineskip}
\section*{Structure and Written Expression}
\subsection*{Sentence Completion}
This section is designed to test your ability to recognize language structures that are appropriate in standard written English. The questions in this section belong to two types, each of which has special directions.\par
DIRECTIONS: Questions 1-15 are partial sentences. Below each sentence you will see four words or phrases, marked (A), (B), (C), and (D). Select the one word or phrase that best completes the sentence. Then, on your answer sheet, find the number of the question and fill in the space that contains the letter for the answer you have chosen. Fill in the space completely.\par
\begin{exampleblock}
\textbf{EXAMPLE I}\par
\noindent Drying flowers is the best way......them.
\begin{enumerate}[label=(\Alph*),leftmargin=*,itemsep=0.2em]
\item to preserve
\item by preserving
\item preserve
\item preserved
\end{enumerate}
\vspace{0.25\baselineskip}
The sentence should state, "Drying flowers is the best way to preserve them." Therefore, the correct answer is (A).\par
\end{exampleblock}
\begin{exampleblock}
\textbf{EXAMPLE II}\par
\noindent Many American universities......as small, private colleges.
\begin{enumerate}[label=(\Alph*),leftmargin=*,itemsep=0.2em]
\item begun
\item beginning
\item began
\item for the beginning
\end{enumerate}
\vspace{0.25\baselineskip}
The sentence should state, "Many American universities began as small, private colleges." Therefore, the correct answer is (C). After you read the directions, begin work on the questions.\par
\end{exampleblock}
\noindent\textbf{1.} The smoke from burning fuels causes pollution if......into the atmosphere.
\begin{enumerate}[label=(\Alph*),leftmargin=*,itemsep=0.2em]
\item it releases
\item it is released
\item it will be released
\item it released
\end{enumerate}
\par
\begingroup
\setlength{\parskip}{0pt}
\setlength{\parindent}{0pt}
\textbf{Review}\par
Answer: B\par
Confidence: high\par
Grammar rule: If-clause passive voice.\par
Klausa if membutuhkan passive voice karena smoke menerima aksi; bentuk yang tepat it is released.\par
\endgroup
\par
\vspace{0.25\baselineskip}
\noindent\textbf{2.} While preparing an issue, newspaper editors decide what......in the editorials.
\begin{enumerate}[label=(\Alph*),leftmargin=*,itemsep=0.2em]
\item viewpoint to take
\item viewpoint takes
\item take a viewpoint
\item takes to a viewpoint
\end{enumerate}
\par
\begingroup
\setlength{\parskip}{0pt}
\setlength{\parindent}{0pt}
\textbf{Review}\par
Answer: A\par
Confidence: high\par
Grammar rule: Noun phrase after an interrogative complement.\par
Setelah decide what, perlu noun phrase + infinitive: viewpoint to take.\par
\endgroup
\par
\vspace{0.25\baselineskip}
\noindent\textbf{3.} Medical researchers claim that each reflex......some stimulus that causes a response.
\begin{enumerate}[label=(\Alph*),leftmargin=*,itemsep=0.2em]
\item involving
\item involvement
\item involves
\item involve
\end{enumerate}
\par
\begingroup
\setlength{\parskip}{0pt}
\setlength{\parindent}{0pt}
\textbf{Review}\par
Answer: C\par
Confidence: high\par
Grammar rule: Subject-verb agreement.\par
Each reflex bersifat tunggal, jadi verb harus singular: involves.\par
\endgroup
\par
\vspace{0.25\baselineskip}
\noindent\textbf{4.} A law of physics stipulates that energy in any system cannot be created or......
\begin{enumerate}[label=(\Alph*),leftmargin=*,itemsep=0.2em]
\item destroy
\item to destroy
\item destroyed
\item destruction
\end{enumerate}
\par
\begingroup
\setlength{\parskip}{0pt}
\setlength{\parindent}{0pt}
\textbf{Review}\par
Answer: C\par
Confidence: high\par
Grammar rule: Parallel passive participles.\par
Parallel dengan be created, gunakan past participle: destroyed.\par
\endgroup
\par
\vspace{0.25\baselineskip}
\noindent\textbf{5.} Until the late 1800s, shopkeepers advertised.......on pictorial signs because their customers were illiterate.
\begin{enumerate}[label=(\Alph*),leftmargin=*,itemsep=0.2em]
\item they produced
\item their products
\item produced their
\item they are produced
\end{enumerate}
\par
\begingroup
\setlength{\parskip}{0pt}
\setlength{\parindent}{0pt}
\textbf{Review}\par
Answer: B\par
Confidence: high\par
Grammar rule: Direct object required.\par
Advertised membutuhkan objek langsung; pilihan yang tepat their products.\par
\endgroup
\par
\vspace{0.25\baselineskip}
\noindent\textbf{6.} Amnesia is the......or total loss of memory concerning past experiences.
\begin{enumerate}[label=(\Alph*),leftmargin=*,itemsep=0.2em]
\item part
\item partially
\item partly
\item partial
\end{enumerate}
\par
\begingroup
\setlength{\parskip}{0pt}
\setlength{\parindent}{0pt}
\textbf{Review}\par
Answer: D\par
Confidence: high\par
Grammar rule: Adjective modifying noun.\par
Setelah the, perlu adjective untuk menerangkan loss: partial.\par
\endgroup
\par
\vspace{0.25\baselineskip}
\noindent\textbf{7.} A sharp sense of smell enables hunting dogs......wild animals.
\begin{enumerate}[label=(\Alph*),leftmargin=*,itemsep=0.2em]
\item to tracking
\item to track
\item to be tracking
\item track them to
\end{enumerate}
\par
\begingroup
\setlength{\parskip}{0pt}
\setlength{\parindent}{0pt}
\textbf{Review}\par
Answer: B\par
Confidence: high\par
Grammar rule: Enable + object + infinitive.\par
Pola enable + object + to + base verb; jadi to track.\par
\endgroup
\par
\vspace{0.25\baselineskip}
\noindent\textbf{8.} The popularity of early melodrama promoted the creation of realistic settings and......effects.
\begin{enumerate}[label=(\Alph*),leftmargin=*,itemsep=0.2em]
\item elaborate special
\item elaborated specially
\item specially elaborating
\item specially to elaborate
\end{enumerate}
\par
\begingroup
\setlength{\parskip}{0pt}
\setlength{\parindent}{0pt}
\textbf{Review}\par
Answer: A\par
Confidence: high\par
Grammar rule: Adjective sequence before noun.\par
Butuh dua adjective sebelum noun effects; bentuk yang tepat elaborate special effects.\par
\endgroup
\par
\vspace{0.25\baselineskip}
\noindent\textbf{9.} The roots anchor a tree in the ground and absorb water......the soil.
\begin{enumerate}[label=(\Alph*),leftmargin=*,itemsep=0.2em]
\item to
\item from
\item into
\item due to
\end{enumerate}
\par
\begingroup
\setlength{\parskip}{0pt}
\setlength{\parindent}{0pt}
\textbf{Review}\par
Answer: B\par
Confidence: high\par
Grammar rule: Preposition choice.\par
Verb absorb diikuti from untuk sumber; absorb water from the soil.\par
\endgroup
\par
\vspace{0.25\baselineskip}
\noindent\textbf{10.} Viruses are so tiny that......only by means of an electron microscope.
\begin{enumerate}[label=(\Alph*),leftmargin=*,itemsep=0.2em]
\item they can see
\item they can be seen
\item can be seen
\item can see them
\end{enumerate}
\par
\begingroup
\setlength{\parskip}{0pt}
\setlength{\parindent}{0pt}
\textbf{Review}\par
Answer: B\par
Confidence: high\par
Grammar rule: Full passive clause after so...that.\par
Setelah so...that perlu klausa lengkap dengan subjek; gunakan they can be seen.\par
\endgroup
\par
\vspace{0.25\baselineskip}
\noindent\textbf{11.} When Arturo Toscanini began his career music was regarded as one......expressing emotions.
\begin{enumerate}[label=(\Alph*),leftmargin=*,itemsep=0.2em]
\item means by
\item meant by it
\item means of
\item the meaning of
\end{enumerate}
\par
\begingroup
\setlength{\parskip}{0pt}
\setlength{\parindent}{0pt}
\textbf{Review}\par
Answer: C\par
Confidence: high\par
Grammar rule: Fixed expression means of.\par
Ekspresi baku one means of + V-ing; jadi means of expressing.\par
\endgroup
\par
\vspace{0.25\baselineskip}
\noindent\textbf{12.} In diving competitions, women perform......men do.
\begin{enumerate}[label=(\Alph*),leftmargin=*,itemsep=0.2em]
\item dive the same as
\item the same dives as
\item dive the same way as
\item the diving is the same
\end{enumerate}
\par
\begingroup
\setlength{\parskip}{0pt}
\setlength{\parindent}{0pt}
\textbf{Review}\par
Answer: B\par
Confidence: high\par
Grammar rule: Object with perform; parallel comparison.\par
Perform membutuhkan objek; bentuk paralel: the same dives as men do.\par
\endgroup
\par
\vspace{0.25\baselineskip}
\noindent\textbf{13.} Syllables are the......of a word according to pronunciation.
\begin{enumerate}[label=(\Alph*),leftmargin=*,itemsep=0.2em]
\item naturally divided
\item divided by nature
\item natural divisions
\item dividing them naturally
\end{enumerate}
\par
\begingroup
\setlength{\parskip}{0pt}
\setlength{\parindent}{0pt}
\textbf{Review}\par
Answer: C\par
Confidence: high\par
Grammar rule: Predicate nominative.\par
Setelah are perlu noun phrase; natural divisions.\par
\endgroup
\par
\vspace{0.25\baselineskip}
\noindent\textbf{14.} Mesopotamian civilizations did not......history until meaning-symbol correspondences had been devised.
\begin{enumerate}[label=(\Alph*),leftmargin=*,itemsep=0.2em]
\item beginning to record
\item record the beginning
\item began recording
\item begin to record
\end{enumerate}
\par
\begingroup
\setlength{\parskip}{0pt}
\setlength{\parindent}{0pt}
\textbf{Review}\par
Answer: D\par
Confidence: high\par
Grammar rule: Did not + base verb infinitive.\par
Setelah did not, gunakan base verb; begin to record.\par
\endgroup
\par
\vspace{0.25\baselineskip}
\noindent\textbf{15.} Before the 1700s, when children worked together with adults, childhood......did not exist.
\begin{enumerate}[label=(\Alph*),leftmargin=*,itemsep=0.2em]
\item as we have known
\item as we know it
\item it is known
\item is known as
\end{enumerate}
\par
\begingroup
\setlength{\parskip}{0pt}
\setlength{\parindent}{0pt}
\textbf{Review}\par
Answer: B\par
Confidence: high\par
Grammar rule: Idiomatic expression.\par
Idiom yang tepat adalah as we know it.\par
\endgroup
\par
\vspace{0.25\baselineskip}
\subsection*{Error Identification}
DIRECTIONS: In questions 16-40 every sentence has four words or phrases that are underlined. The four underlined portions of each sentence are marked (A), (B), (C), and (D). Identify the one word or phrase that makes the sentence incorrect. Then, on your answer sheet, find the number of the question and fill in the space that contains the letter for the answer you have chosen. Fill in the space completely.\par
\begin{exampleblock}
\textbf{EXAMPLE I}\par
Christopher Columbus \uline{has sailed}(A) from Europe in 1492 and \uline{discovered a}(B) \uline{new land}(C) he \uline{thought to}(D) be India.\par
The sentence should state, "Christopher Columbus sailed from Europe in 1492 and discovered a new land he thought to be India." Therefore, you should choose answer (A).\par
\end{exampleblock}
\begin{exampleblock}
\textbf{EXAMPLE II}\par
\uline{As the roles}(A) of people in society change, \uline{so does}(B) the \uline{rules of}(C) conduct in certain \uline{situations}(D).\par
The sentence should state, "As the roles of people in society change, so do the rules of\par
conduct in certain situations." Therefore, you should choose answer (B).\par
\end{exampleblock}
\noindent\textbf{16.} Carl Anderson \uline{discovered}(A) two atomic \uline{particles}(B) that he \uline{identified}(C) while \uline{studied}(D) cosmic rays.
\par
\begingroup
\setlength{\parskip}{0pt}
\setlength{\parindent}{0pt}
\textbf{Review}\par
Answer: D\par
Confidence: high\par
Grammar rule: After while, use participle/gerund.\par
Correct form: studying\par
Setelah while gunakan bentuk V-ing; bagian (D) studied harus studying.\par
\endgroup
\par
\vspace{\questionblockgap}
\noindent\textbf{17.} No one \uline{knows}(A) exactly how many \uline{species}(B) of \uline{animals}(C) \uline{lives}(D) on earth.
\par
\begingroup
\setlength{\parskip}{0pt}
\setlength{\parindent}{0pt}
\textbf{Review}\par
Answer: D\par
Confidence: high\par
Grammar rule: Plural subject takes plural verb.\par
Correct form: live\par
Subjek species of animals bersifat jamak; verb harus live, jadi (D) salah.\par
\endgroup
\par
\vspace{\questionblockgap}
\noindent\textbf{18.} Assessment instruments in nursery schools \uline{they feature}(A) items and \uline{other materials}(B) different \uline{from those}(C) on elementary \uline{school tests}(D).
\par
\begingroup
\setlength{\parskip}{0pt}
\setlength{\parindent}{0pt}
\textbf{Review}\par
Answer: A\par
Confidence: high\par
Grammar rule: Double-subject error.\par
Correct form: feature\par
Kalimat sudah punya subjek Assessment instruments; they tidak perlu. (A) salah, gunakan feature.\par
\endgroup
\par
\vspace{\questionblockgap}
\noindent\textbf{19.} Michigan's rivers, inlets,\uline{ and lakes}(A) attract tourists \uline{who derive}(B) pleasure \uline{from canoeing}(C) and \uline{water-ski}(D).
\par
\begingroup
\setlength{\parskip}{0pt}
\setlength{\parindent}{0pt}
\textbf{Review}\par
Answer: D\par
Confidence: high\par
Grammar rule: Parallel gerund after from.\par
Correct form: water-skiing\par
Setelah from, dua aktivitas harus parallel gerund; water-ski -> water-skiing (D).\par
\endgroup
\par
\vspace{\questionblockgap}
\noindent\textbf{20.} Analysts \uline{have translated}(A) clay \uline{tables that}(B) demonstrate \uline{that the}(C) Babylonians were \uline{high skilled}(D) in arithmetic.
\par
\begingroup
\setlength{\parskip}{0pt}
\setlength{\parindent}{0pt}
\textbf{Review}\par
Answer: D\par
Confidence: high\par
Grammar rule: Adverb modifying adjective.\par
Correct form: highly skilled\par
Adverb dibutuhkan untuk memodifikasi skilled; seharusnya highly skilled (D).\par
\endgroup
\par
\vspace{\questionblockgap}
\noindent\textbf{21.} The visual \uline{nerves of}(A) the brain \uline{interprets}(B) \uline{wavelengths}(C) of \uline{light as}(D) perceptions of color.
\par
\begingroup
\setlength{\parskip}{0pt}
\setlength{\parindent}{0pt}
\textbf{Review}\par
Answer: B\par
Confidence: high\par
Grammar rule: Subject-verb agreement.\par
Correct form: interpret\par
Subjek plural nerves membutuhkan verb plural; interpret, bukan interprets (B).\par
\endgroup
\par
\vspace{\questionblockgap}
\noindent\textbf{22.} \uline{It is}(A) possible to have wealth but \uline{little()B}() income and \uline{having}(C) income but \uline{no wealth}(D).
\par
\begingroup
\setlength{\parskip}{0pt}
\setlength{\parindent}{0pt}
\textbf{Review}\par
Answer: C\par
Confidence: medium\par
Grammar rule: Parallel infinitive structure.\par
Correct form: to have\par
Struktur paralel setelah possible to have; harus to have ... and to have ... Jadi (C) salah.\par
\endgroup
\par
\vspace{\questionblockgap}
\noindent\textbf{23.} When a criminal case goes to \uline{trial}(A), the defendant may \uline{election}(B) to have it heard \uline{either}(C) by a jury or by \uline{a judge}(D).
\par
\begingroup
\setlength{\parskip}{0pt}
\setlength{\parindent}{0pt}
\textbf{Review}\par
Answer: B\par
Confidence: high\par
Grammar rule: Modal + base verb.\par
Correct form: elect\par
Modal + base verb; election harus elect (B).\par
\endgroup
\par
\vspace{\questionblockgap}
\noindent\textbf{24.} John Keynes \uline{used}(A) his \uline{knowledges}(B) of economics to \uline{help}(C) his college and \uline{himself}(D).
\par
\begingroup
\setlength{\parskip}{0pt}
\setlength{\parindent}{0pt}
\textbf{Review}\par
Answer: B\par
Confidence: high\par
Grammar rule: Uncountable noun.\par
Correct form: knowledge\par
Knowledge uncountable; tidak pakai -s. (B) salah.\par
\endgroup
\par
\vspace{\questionblockgap}
\noindent\textbf{25.} Government \uline{offices store}(A) and maintain \uline{such documents}(B) as certificates of birth, \uline{married}(C), and \uline{death}(D).
\par
\begingroup
\setlength{\parskip}{0pt}
\setlength{\parindent}{0pt}
\textbf{Review}\par
Answer: C\par
Confidence: high\par
Grammar rule: Noun parallelism after certificates of.\par
Correct form: marriage\par
Parallel noun setelah certificates of; married harus marriage (C).\par
\endgroup
\par
\vspace{\questionblockgap}
\noindent\textbf{26.} \uline{After the}(A) Constitution was \uline{signed}(B), Delaware became \uline{the first}(C) state to \uline{ratifying}(D) it.
\par
\begingroup
\setlength{\parskip}{0pt}
\setlength{\parindent}{0pt}
\textbf{Review}\par
Answer: D\par
Confidence: high\par
Grammar rule: To + base verb.\par
Correct form: ratify\par
to + base verb; ratifying harus ratify (D).\par
\endgroup
\par
\vspace{\questionblockgap}
\noindent\textbf{27.} Migrant \uline{workers live}(A) in \uline{substandard}(B) \uline{unsanitary}(C), and dilapidated housing and often \uline{are lacking}(D) medical care.
\par
\begingroup
\setlength{\parskip}{0pt}
\setlength{\parindent}{0pt}
\textbf{Review}\par
Answer: D\par
Confidence: high\par
Grammar rule: Simple active verb needed.\par
Correct form: lack\par
Butuh verb aktif sederhana; gunakan lack, bukan are lacking (D).\par
\endgroup
\par
\vspace{\questionblockgap}
\noindent\textbf{28.} The \uline{mining}(A) of minerals often \uline{bring about}(B) \uline{the destruction}(C) of landscapes and \uline{wildlife}(D) habitats.
\par
\begingroup
\setlength{\parskip}{0pt}
\setlength{\parindent}{0pt}
\textbf{Review}\par
Answer: B\par
Confidence: high\par
Grammar rule: Singular subject with singular verb.\par
Correct form: brings about\par
Subjek tunggal mining -> verb singular brings about (B).\par
\endgroup
\par
\vspace{\questionblockgap}
\noindent\textbf{29.} Christopher Marlowe \uline{established}(A) his \uline{theatrical}(B) reputation with Tamburlaine the Great, \uline{written}(C) in high verse and \uline{reflected}(D) his unconventional thought.
\par
\begingroup
\setlength{\parskip}{0pt}
\setlength{\parindent}{0pt}
\textbf{Review}\par
Answer: D\par
Confidence: medium\par
Grammar rule: Parallel reduced modifiers.\par
Correct form: reflecting\par
Parallel dengan written (participle); gunakan reflecting, bukan reflected (D).\par
\endgroup
\par
\vspace{\questionblockgap}
\noindent\textbf{30.} William H . Bonney, \uline{better known}(A) as Billy the Kid, \uline{shoot a man}(B) to \uline{death in a quarrel}(C) and \uline{had to flee}(D) to New Mexico.
\par
\begingroup
\setlength{\parskip}{0pt}
\setlength{\parindent}{0pt}
\textbf{Review}\par
Answer: B\par
Confidence: high\par
Grammar rule: Past tense needed.\par
Correct form: shot a man\par
Narasi lampau; shoot harus shot (B).\par
\endgroup
\par
\vspace{\questionblockgap}
\noindent\textbf{31.} Foxes stay in \uline{closely knit}(A) family groups \uline{while the}(B) \uline{young ones}(C) are \uline{grow up}(D).
\par
\begingroup
\setlength{\parskip}{0pt}
\setlength{\parindent}{0pt}
\textbf{Review}\par
Answer: D\par
Confidence: high\par
Grammar rule: Progressive form after are.\par
Correct form: growing up\par
Setelah are gunakan present participle; are growing up (D).\par
\endgroup
\par
\vspace{\questionblockgap}
\noindent\textbf{32.} By 1938, the \uline{sale}(A) of \uline{records album}(B) \uline{had reached}(C) \$26 \uline{million}(D) a year.
\par
\begingroup
\setlength{\parskip}{0pt}
\setlength{\parindent}{0pt}
\textbf{Review}\par
Answer: B\par
Confidence: medium\par
Grammar rule: Compound noun form.\par
Correct form: record albums\par
Compound noun plural; records album -> record albums (B).\par
\endgroup
\par
\vspace{\questionblockgap}
\noindent\textbf{33.} Egyptian artisans \uline{made}(A) glass that was \uline{colored}(B) by the \uline{present}(C) of \uline{impurities}(D).
\par
\begingroup
\setlength{\parskip}{0pt}
\setlength{\parindent}{0pt}
\textbf{Review}\par
Answer: C\par
Confidence: high\par
Grammar rule: Noun needed after the.\par
Correct form: presence\par
Noun yang benar setelah the; present harus presence (C).\par
\endgroup
\par
\vspace{\questionblockgap}
\noindent\textbf{34.} A theory \uline{called}(A) plate tectonics \uline{explain}(B) the \uline{formation}(C) of the surface \uline{features}(D) of the earth.
\par
\begingroup
\setlength{\parskip}{0pt}
\setlength{\parindent}{0pt}
\textbf{Review}\par
Answer: B\par
Confidence: high\par
Grammar rule: Singular subject-verb agreement.\par
Correct form: explains\par
Subject singular theory -> verb explains (B).\par
\endgroup
\par
\vspace{\questionblockgap}
\noindent\textbf{35.} Members of \uline{high school}(A) clubs learn to \uline{participation}(B) in teams through their \uline{involvement}(C) in \uline{community projects}(D).
\par
\begingroup
\setlength{\parskip}{0pt}
\setlength{\parindent}{0pt}
\textbf{Review}\par
Answer: B\par
Confidence: high\par
Grammar rule: Learn to + base verb.\par
Correct form: participate\par
learn to + base verb; participation -> participate (B).\par
\endgroup
\par
\vspace{\questionblockgap}
\noindent\textbf{36.} When \uline{too many}(A) firms enter \uline{competitive}(B) markets, their share \uline{of profits}(C) \uline{will fell}(D).
\par
\begingroup
\setlength{\parskip}{0pt}
\setlength{\parindent}{0pt}
\textbf{Review}\par
Answer: D\par
Confidence: high\par
Grammar rule: Will + base verb.\par
Correct form: fall\par
will + base verb; fell -> fall (D).\par
\endgroup
\par
\vspace{\questionblockgap}
\noindent\textbf{37.} The term "middle class" \uline{describes}(A) people \uline{between}(B) the upper and \uline{the low}(C) \uline{social}(D) classes.
\par
\begingroup
\setlength{\parskip}{0pt}
\setlength{\parindent}{0pt}
\textbf{Review}\par
Answer: C\par
Confidence: high\par
Grammar rule: Parallel adjective form.\par
Correct form: lower\par
Parallel adjective upper-lower; low harus lower (C).\par
\endgroup
\par
\vspace{\questionblockgap}
\noindent\textbf{38.} Copper comes from seven \uline{types}(A) of \uline{ores}(B) that \uline{also contain}(C) \uline{the other}(D) materials.
\par
\begingroup
\setlength{\parskip}{0pt}
\setlength{\parindent}{0pt}
\textbf{Review}\par
Answer: D\par
Confidence: high\par
Grammar rule: No article before plural noun here.\par
Correct form: other\par
Other + plural noun tanpa the; jadi other materials (D).\par
\endgroup
\par
\vspace{\questionblockgap}
\noindent\textbf{39.} Matthew Henson \uline{received}(A) many \uline{honor}(B) for \uline{his part}(C) in \uline{the expedition}(D) to the North Pole.
\par
\begingroup
\setlength{\parskip}{0pt}
\setlength{\parindent}{0pt}
\textbf{Review}\par
Answer: B\par
Confidence: high\par
Grammar rule: Many + plural count noun.\par
Correct form: honors\par
Many + plural count noun; honor -> honors (B).\par
\endgroup
\par
\vspace{\questionblockgap}
\noindent\textbf{40.} Silicon chips contain \uline{thousands}(A) of \uline{circuits}(B) in an area \uline{as smaller}(C) than \uline{a fingernail}(D).
\par
\begingroup
\setlength{\parskip}{0pt}
\setlength{\parindent}{0pt}
\textbf{Review}\par
Answer: C\par
Confidence: high\par
Grammar rule: Comparative structure with than.\par
Correct form: smaller\par
Comparative yang tepat smaller than; hapus as, jadi (C) salah.\par
\endgroup
\par
\vspace{\questionblockgap}
\section*{Reading Comprehension}
\begin{center}
\textbf{Time 55 minutes}
\end{center}
DIRECTIONS: In this section you will read several passages. Each passage is followed by a series of questions. For questions 1-50, you need to select the best answer, (A), (B), (C), or (D), to each question. Then, on your answer sheet, find the number of the question and fill in the space that contains the letter of the answer you have selected. Fill in the space completely. Answer all questions following a passage on the basis of what is stated or implied in the passage.\par
\begin{exampleblock}
\small
\sloppy
Read the following passage:\par
A tomahawk is a small ax used as a tool and a weapon by the North American Indian\par
tribes. An average tomahawk was not very long and did not weigh a great deal. Originally,\par
the head of the tomahawk was made of a shaped stone or an animal bone and was mounted on\par
Line a wooden handle. After the arrival of the European settlers, the Indians began to use toma-\par
hawks with iron heads. Indian males and females of all ages used tomahawks to chop and cut wood, pound stakes into the ground to put up wigwams, and do many other chores. Indian warriors relied on tomahawks as weapons and even threw them at their enemies. Some types of tomahawks were used in religious ceremonies. Contemporary American idioms reflect\par
this aspect of American heritage.\par
\end{exampleblock}
\begin{exampleblock}
\textbf{EXAMPLE I}\par
\noindent Early tomahawk heads were made of
\begin{enumerate}[label=(\Alph*),leftmargin=*,itemsep=0.2em]
\item stone or bone
\item wood or sticks
\item European iron
\item religious weapons
\end{enumerate}
\vspace{0.25\baselineskip}
According to the passage, early tomahawk heads were made of stone or bone. Therefore, the correct answer is (A).\par
\end{exampleblock}
\begin{exampleblock}
\textbf{EXAMPLE II}\par
\noindent How has the Indian use of tomahawks affected American daily life today?
\begin{enumerate}[label=(\Alph*),leftmargin=*,itemsep=0.2em]
\item Tomahawks are still used as weapons.
\item Tomahawks are used as tools for certain jobs.
\item Contemporary language refers to tomahawks.
\item Indian tribes cherish tomahawks as heirlooms.
\end{enumerate}
\vspace{0.25\baselineskip}
The passage states, "Contemporary American idioms reflect this aspect of American heritage." The correct answer is (C). After you read the directions, begin work on the questions.\par
\end{exampleblock}
\subsection*{Questions 1-12}
\begingroup
\small
\sloppy
\setlength{\parskip}{0pt}
\setlength{\parindent}{0pt}
Even with his diverse experience as an elected official at the state level, Andrew Johnson\par
was the first president of the United States ever to be impeached, primarily because of his\par
violent temper and unyielding stubbornness. His career started in 1828 with his election to\par
the city council of Greenville, Tennessee, and after two years as an alderman, he took office\par
as mayor. His advancements followed in rapid succession when he was elected to the Ten-\par
nessee state senate, then as the state governor, and later to the U.S. House of Representatives\par
for five consecutive terms.\par
In 1864, Johnson ran for the office of vice-president on the Lincoln-Johnson ticket and\par
was inaugurated in 1865. After Lincoln's assassination six weeks into his term, Johnson\par
found himself president at a time when southern leaders were concerned about their forced\par
alliance with the northern states and feared retaliation for their support of the secession. In-\par
stead, however, with the diplomatic skill he had learned from Lincoln, Johnson offered full\par
pardon to almost all Confederates on the condition that they take an oath of allegiance. He\par
further reorganized the former Confederate states and set up legislative elections.\par
Congressional opposition to his peace-making policies resulted in gridlock between the\par
House and Johnson, and the stalemate grew into an open conflict on the issue of the emanci-\par
pation of slaves. While Johnson held the view that newly freed slaves lacked understanding\par
and knowledge of civil liberties to vote intelligently, Congress overrode Johnson's veto of\par
the Civil Rights Bill, which awarded them citizenship and ratified the Fourteenth Amend-\par
ment. In the years that followed, Congress passed bills depriving the president of the power\par
to pardon political criminals, stripping away his status of commander-in-chief, and taking\par
away Johnson's right to dismiss civil and executive officers from their duties. Johnson ve-\par
toed each bill, and each veto was overridden. When Johnson dismissed the secretary of war,\par
Edwin Stanton, Stanton refused to step down and was supported by the House of Repre-\par
sentatives, which voted to impeach Johnson. At the trial, the Senate came one vote short of\par
the two-thirds majority necessary to remove him from office. After Johnson's term expired,\par
he returned to his home state, but in 1875 he was elected senator and went back to Wash-\par
ington to take his seat.\par
\endgroup
\medskip
\noindent\textbf{1.} What does the passage mainly discuss?
\begin{enumerate}[label=(\Alph*),leftmargin=*,itemsep=0.2em]
\item Andrew Johnson's personal characteristics
\item Andrew Johnson's career as a politician
\item Congressional decisions in the late 1800s
\item Congressional decisions and procedures in the late 1800s
\end{enumerate}
\par
\begingroup
\setlength{\parskip}{0pt}
\setlength{\parindent}{0pt}
\textbf{Review}\par
Answer: B\par
Confidence: high\par
Passage ini membahas perjalanan politik Andrew Johnson, masa kepresidenannya, dan konflik impeachment yang ia hadapi.\par
\endgroup
\par
\vspace{0.25\baselineskip}
\noindent\textbf{2.} In line 4, the phrase "took office" is closest in meaning to
\begin{enumerate}[label=(\Alph*),leftmargin=*,itemsep=0.2em]
\item moved into an office
\item became an official
\item began a government job
\item rearranged the office
\end{enumerate}
\par
\begingroup
\setlength{\parskip}{0pt}
\setlength{\parindent}{0pt}
\textbf{Review}\par
Answer: C\par
Confidence: high\par
Evidence refs: line 4\par
Evidence quote: "took office"\par
Di sini, took office berarti mulai menjabat dalam posisi pemerintahan.\par
\endgroup
\par
\vspace{0.25\baselineskip}
\noindent\textbf{3.} What can be inferred from the first paragraph about Andrew Johnson's work in Tennessee?
\begin{enumerate}[label=(\Alph*),leftmargin=*,itemsep=0.2em]
\item His personality precluded him from important positions.
\item His work became known to the governor.
\item He was elected to several important posts.
\item He was represented to the posts five times.
\end{enumerate}
\par
\begingroup
\setlength{\parskip}{0pt}
\setlength{\parindent}{0pt}
\textbf{Review}\par
Answer: C\par
Confidence: high\par
Paragraf pertama menunjukkan bahwa Johnson terpilih ke beberapa jabatan penting di Tennessee.\par
\endgroup
\par
\vspace{0.25\baselineskip}
\noindent\textbf{4.} In line 11, the word "alliance" is closest in meaning to
\begin{enumerate}[label=(\Alph*),leftmargin=*,itemsep=0.2em]
\item union
\item counsel
\item allowance
\item allotment
\end{enumerate}
\par
\begingroup
\setlength{\parskip}{0pt}
\setlength{\parindent}{0pt}
\textbf{Review}\par
Answer: A\par
Confidence: high\par
Evidence refs: line 11\par
Evidence quote: "alliance with the northern states"\par
Alliance di sini berarti persatuan.\par
\endgroup
\par
\vspace{0.25\baselineskip}
\noindent\textbf{5.} According to the passage, what led to Johnson's downfall?
\begin{enumerate}[label=(\Alph*),leftmargin=*,itemsep=0.2em]
\item The state of the nation's economy
\item His liberal position on slavery
\item His personal characteristics
\item His waffling and hesitation
\end{enumerate}
\par
\begingroup
\setlength{\parskip}{0pt}
\setlength{\parindent}{0pt}
\textbf{Review}\par
Answer: C\par
Confidence: high\par
Evidence refs: lines 1-3\par
Evidence quote: "violent temper and unyielding stubbornness"\par
Passage ini mengaitkan kejatuhannya langsung dengan karakter pribadinya.\par
\endgroup
\par
\vspace{0.25\baselineskip}
\noindent\textbf{6.} The author of the passage implies that when Johnson became president he
\begin{enumerate}[label=(\Alph*),leftmargin=*,itemsep=0.2em]
\item was a dedicated supporter of civil rights
\item was a soft-spoken and careful diplomat
\item had an extensive background in politics
\item had already experienced political turmoil
\end{enumerate}
\par
\begingroup
\setlength{\parskip}{0pt}
\setlength{\parindent}{0pt}
\textbf{Review}\par
Answer: C\par
Confidence: high\par
Evidence refs: lines 3-9\par
Evidence quote: "city council ... state senate ... state governor ... U.S. House"\par
Ia sudah memiliki pengalaman politik yang luas sebelum menjadi presiden.\par
\endgroup
\par
\vspace{0.25\baselineskip}
\noindent\textbf{7.} According to the passage, at the beginning of Johnson's term as president southerners were
\begin{enumerate}[label=(\Alph*),leftmargin=*,itemsep=0.2em]
\item expected to secede from the union
\item apprehensive about their future
\item singled out as scapegoats
\item afraid of his violent temper
\end{enumerate}
\par
\begingroup
\setlength{\parskip}{0pt}
\setlength{\parindent}{0pt}
\textbf{Review}\par
Answer: B\par
Confidence: high\par
Evidence refs: lines 10-11\par
Evidence quote: "concerned ... feared retaliation"\par
Detail ini menunjukkan bahwa warga Selatan cemas terhadap masa depan mereka.\par
\endgroup
\par
\vspace{0.25\baselineskip}
\noindent\textbf{8.} According to the passage, Congress's disapproval of Andrew Johnson's policies was
\begin{enumerate}[label=(\Alph*),leftmargin=*,itemsep=0.2em]
\item short-lived and groundless
\item detrimental to his presidency
\item directed at his civic duties
\item stopped as soon as it emerged
\end{enumerate}
\par
\begingroup
\setlength{\parskip}{0pt}
\setlength{\parindent}{0pt}
\textbf{Review}\par
Answer: B\par
Confidence: high\par
Evidence refs: lines 15-25\par
Evidence quote: "gridlock ... veto was overridden ... voted to impeach Johnson"\par
Oposisi di Kongres merusak dan melemahkan masa kepresidenannya.\par
\endgroup
\par
\vspace{0.25\baselineskip}
\noindent\textbf{9.} In line 21, the word "pardon" is closest in meaning to
\begin{enumerate}[label=(\Alph*),leftmargin=*,itemsep=0.2em]
\item parade
\item patronize
\item exonerate
\item extricate
\end{enumerate}
\par
\begingroup
\setlength{\parskip}{0pt}
\setlength{\parindent}{0pt}
\textbf{Review}\par
Answer: C\par
Confidence: high\par
Evidence refs: line 21\par
Evidence quote: "pardon political criminals"\par
Pardon paling dekat maknanya dengan exonerate dalam konteks ini.\par
\endgroup
\par
\vspace{0.25\baselineskip}
\noindent\textbf{10.} The author of the passage implies that the Stanton affair proved the president's
\begin{enumerate}[label=(\Alph*),leftmargin=*,itemsep=0.2em]
\item lack of stamina
\item lack of electoral vote
\item loss of willpower
\item loss of authority
\end{enumerate}
\par
\begingroup
\setlength{\parskip}{0pt}
\setlength{\parindent}{0pt}
\textbf{Review}\par
Answer: D\par
Confidence: high\par
Evidence refs: lines 22-25\par
Evidence quote: "Stanton refused to step down"\par
Peristiwa Stanton menunjukkan bahwa Johnson kehilangan otoritas eksekutif.\par
\endgroup
\par
\vspace{0.25\baselineskip}
\noindent\textbf{11.} In line 23, the word "dismissed" is closest in meaning to
\begin{enumerate}[label=(\Alph*),leftmargin=*,itemsep=0.2em]
\item distanced
\item fired
\item disdained
\item flounced
\end{enumerate}
\par
\begingroup
\setlength{\parskip}{0pt}
\setlength{\parindent}{0pt}
\textbf{Review}\par
Answer: B\par
Confidence: high\par
Evidence refs: line 23\par
Evidence quote: "Johnson dismissed the secretary of war"\par
Dismissed berarti diberhentikan atau dicopot dari jabatan.\par
\endgroup
\par
\vspace{0.25\baselineskip}
\noindent\textbf{12.} According to the passage, the attempt to impeach Andrew Johnson
\begin{enumerate}[label=(\Alph*),leftmargin=*,itemsep=0.2em]
\item succeeded as expected by the House
\item failed by a minimal margin
\item put an end to his political career
\item overwhelmed his supporters in Tennessee
\end{enumerate}
\par
\begingroup
\setlength{\parskip}{0pt}
\setlength{\parindent}{0pt}
\textbf{Review}\par
Answer: B\par
Confidence: high\par
Evidence refs: lines 25-26\par
Evidence quote: "the Senate came one vote short"\par
Upaya impeachment gagal dengan selisih yang sangat tipis.\par
\endgroup
\par
\vspace{0.25\baselineskip}
\subsection*{Questions 13-24}
\begingroup
\small
\sloppy
\setlength{\parskip}{0pt}
\setlength{\parindent}{0pt}
Sex-trait stereotypes may be defined as a set of psychological attributes that characterize\par
men more frequently than women. Thus, males are often described as ambitious, unemo-\par
tional, and independent and, on the other hand, selfish, unrefined, and insensitive. Females\par
are described as emotional, irrational, high-strung, and tentative. In spite of the egalitarian\par
movement, recent studies have demonstrated that sex-trait stereotypes remain common\par
among young adults today. In fact, such stereotyping has proved to be the psychological\par
justification for social beliefs concerning the appropriateness of various activities for men\par
and women that further perpetuate the different sex roles traditionally ascribed to men and\par
women.\par
The awareness of sex-trait stereotypes in the United States develops in a linear fashion be-\par
tween the ages of four and ten. Generally, knowledge of male stereotypical characteristics\par
develops earlier, whereas knowledge of female characteristics increases more rapidly be-\par
tween the ages of four and seven. While the reasons for this learning are not fully under-\par
stood, evidence suggests that at the preschool level children's literature and television pro-\par
grams provide powerful models and reinforcement for stereotyped views.\par
Studies designed to compare sex-trait stereotypes cross-nationally show a high degree of\par
correspondence in the characteristics ascribed to men and women. As findings have been\par
obtained in other countries, two hypotheses have been advanced to explain the commonali-\par
ties in sex trait stereotyping. One states that pancultural similarities play a role in the psy-\par
chological characteristics attributed to men and women, and the second states that the gen-\par
eral picture is one of cultural relativism.\par
\endgroup
\medskip
\noindent\textbf{13.} Which of the following is the best title for the passage?
\begin{enumerate}[label=(\Alph*),leftmargin=*,itemsep=0.2em]
\item A Relativist Perspective on Stereotyping
\item The Pervasiveness of Sex-Trait Stereotypes
\item A Unilateral Approach to Sex-Trait Stereotyping
\item A Cross-examination of Stereotypical Behaviors
\end{enumerate}
\par
\begingroup
\setlength{\parskip}{0pt}
\setlength{\parindent}{0pt}
\textbf{Review}\par
Answer: B\par
Confidence: high\par
Passage ini membahas betapa luasnya stereotip sex-trait serta bagaimana stereotip itu dipelajari dan disebarkan.\par
\endgroup
\par
\vspace{0.25\baselineskip}
\noindent\textbf{14.} In line 2, the word "ambitious" is closest in meaning to
\begin{enumerate}[label=(\Alph*),leftmargin=*,itemsep=0.2em]
\item enterprising
\item ambiguous
\item anxious
\item honest
\end{enumerate}
\par
\begingroup
\setlength{\parskip}{0pt}
\setlength{\parindent}{0pt}
\textbf{Review}\par
Answer: A\par
Confidence: high\par
Evidence refs: line 2\par
Evidence quote: "ambitious"\par
Ambitious paling dekat maknanya dengan enterprising.\par
\endgroup
\par
\vspace{0.25\baselineskip}
\noindent\textbf{15.} In line 4, the word "high-strung" is closest in meaning to
\begin{enumerate}[label=(\Alph*),leftmargin=*,itemsep=0.2em]
\item high-class
\item fair-minded
\item nervous
\item hideous
\end{enumerate}
\par
\begingroup
\setlength{\parskip}{0pt}
\setlength{\parindent}{0pt}
\textbf{Review}\par
Answer: C\par
Confidence: high\par
Evidence refs: line 4\par
Evidence quote: "high-strung"\par
High-strung berarti nervous atau tense.\par
\endgroup
\par
\vspace{0.25\baselineskip}
\noindent\textbf{16.} Which of the following statements is supported in the passage?
\begin{enumerate}[label=(\Alph*),leftmargin=*,itemsep=0.2em]
\item The egalitarian movement has been a resounding success.
\item The beliefs of young adults have shown little change.
\item Young adults have participated in many common studies.
\item The beliefs of young adults are more common among the old.
\end{enumerate}
\par
\begingroup
\setlength{\parskip}{0pt}
\setlength{\parindent}{0pt}
\textbf{Review}\par
Answer: B\par
Confidence: high\par
Evidence refs: lines 5-6\par
Evidence quote: "remain common among young adults today"\par
Passage ini menyatakan bahwa keyakinan kaum muda hampir tidak berubah.\par
\endgroup
\par
\vspace{0.25\baselineskip}
\noindent\textbf{17.} In line 8, the word "perpetuate" is closest in meaning to
\begin{enumerate}[label=(\Alph*),leftmargin=*,itemsep=0.2em]
\item personalize
\item perplex
\item maintain
\item mount
\end{enumerate}
\par
\begingroup
\setlength{\parskip}{0pt}
\setlength{\parindent}{0pt}
\textbf{Review}\par
Answer: C\par
Confidence: high\par
Evidence refs: line 8\par
Evidence quote: "perpetuate"\par
Perpetuate berarti mempertahankan atau melanjutkan.\par
\endgroup
\par
\vspace{0.25\baselineskip}
\noindent\textbf{18.} It can be inferred from the passage that social beliefs precipitate
\begin{enumerate}[label=(\Alph*),leftmargin=*,itemsep=0.2em]
\item the on-going egalitarian change
\item the rationalization for stereotyping
\item nontraditional gender roles
\item concerns for the legitimacy of sex traits
\end{enumerate}
\par
\begingroup
\setlength{\parskip}{0pt}
\setlength{\parindent}{0pt}
\textbf{Review}\par
Answer: B\par
Confidence: high\par
Pilihan B sesuai dengan pernyataan bahwa keyakinan sosial memicu rasionalisasi stereotyping.\par
\endgroup
\par
\vspace{0.25\baselineskip}
\noindent\textbf{19.} It can be inferred from the second paragraph that young children learn about sex-trait stereotypes
\begin{enumerate}[label=(\Alph*),leftmargin=*,itemsep=0.2em]
\item by watching their parents
\item by being exposed to various media
\item after they start school
\item when their learning is reinforced
\end{enumerate}
\par
\begingroup
\setlength{\parskip}{0pt}
\setlength{\parindent}{0pt}
\textbf{Review}\par
Answer: B\par
Confidence: high\par
Evidence refs: lines 13-15\par
Evidence quote: "children's literature and television programs provide powerful models"\par
Paragraf kedua mendukung gagasan bahwa pembelajaran terjadi melalui paparan berbagai media.\par
\endgroup
\par
\vspace{0.25\baselineskip}
\noindent\textbf{20.} Where in the passage does the author refer to limitations of sex-trait research?
\begin{enumerate}[label=(\Alph*),leftmargin=*,itemsep=0.2em]
\item Lines 1-3
\item Lines 4-6
\item Lines 13-15
\item Lines 16-19
\end{enumerate}
\par
\begingroup
\setlength{\parskip}{0pt}
\setlength{\parindent}{0pt}
\textbf{Review}\par
Answer: C\par
Confidence: high\par
Evidence refs: lines 13-15\par
Evidence quote: "the reasons for this learning are not fully understood"\par
Batasan penelitian disebut pada baris 13-15.\par
\endgroup
\par
\vspace{0.25\baselineskip}
\noindent\textbf{21.} According to the passage, characterizations of men and women as having particular sets of attributes are
\begin{enumerate}[label=(\Alph*),leftmargin=*,itemsep=0.2em]
\item reflected in modern fashion
\item found in several countries
\item uniform across all groups
\item contingent on a socioeconomic class
\end{enumerate}
\par
\begingroup
\setlength{\parskip}{0pt}
\setlength{\parindent}{0pt}
\textbf{Review}\par
Answer: B\par
Confidence: high\par
Evidence refs: lines 16-18\par
Evidence quote: "cross-nationally show a high degree of correspondence"\par
Passage ini mengatakan bahwa karakterisasi tersebut ditemukan di beberapa negara.\par
\endgroup
\par
\vspace{0.25\baselineskip}
\noindent\textbf{22.} In line 17, the word "correspondence" is closest in meaning to
\begin{enumerate}[label=(\Alph*),leftmargin=*,itemsep=0.2em]
\item letters
\item writing
\item agreement
\item discord
\end{enumerate}
\par
\begingroup
\setlength{\parskip}{0pt}
\setlength{\parindent}{0pt}
\textbf{Review}\par
Answer: C\par
Confidence: high\par
Evidence refs: line 17\par
Evidence quote: "correspondence"\par
Correspondence paling dekat maknanya dengan agreement.\par
\endgroup
\par
\vspace{0.25\baselineskip}
\noindent\textbf{23.} The author of the passage would most probably agree with which of the following statements?
\begin{enumerate}[label=(\Alph*),leftmargin=*,itemsep=0.2em]
\item Social attitudes toward women have been updated and made more balanced.
\item Social attitudes toward men are continually nullified and modernized.
\item The women's liberation movement has borne little fruit.
\item Social attitudes are not likely to change radically.
\end{enumerate}
\par
\begingroup
\setlength{\parskip}{0pt}
\setlength{\parindent}{0pt}
\textbf{Review}\par
Answer: D\par
Confidence: medium\par
Karena stereotip tetap umum meski ada gerakan reformasi, penulis memandang perubahan radikal tidak mungkin terjadi.\par
\endgroup
\par
\vspace{0.25\baselineskip}
\noindent\textbf{24.} The passage is probably an excerpt from an article on
\begin{enumerate}[label=(\Alph*),leftmargin=*,itemsep=0.2em]
\item demographics
\item sociology
\item sociobiology
\item psychotherapy
\end{enumerate}
\par
\begingroup
\setlength{\parskip}{0pt}
\setlength{\parindent}{0pt}
\textbf{Review}\par
Answer: B\par
Confidence: high\par
Topiknya adalah stereotyping, peran sosial, dan sikap sosial, sehingga cocok dengan bidang sociology.\par
\endgroup
\par
\vspace{0.25\baselineskip}
\subsection*{Questions 25-32}
\begingroup
\small
\sloppy
\setlength{\parskip}{0pt}
\setlength{\parindent}{0pt}
There are many reasons why food fads have continued to flourish. Garlic has long been\par
touted as an essential ingredient of physical prowess and as a flu remedy, squash has been\par
thought by some to cure digestive disorders, and red pepper has been alleged to promote en-\par
durance. The natural human desire for a simple solution to a difficult problem sets the stage\par
for promoting miraculous potions, pills, and combinations of chemicals. The gullible indi-\par
viduals who eagerly embrace any second-hand information with scientific overtones pro-\par
vide the foundation for healthy business enterprises.\par
A person who has never crossed the threshold of a health food store may be astounded, be-\par
wildered, or overjoyed. Countless elixirs, herbs, powders, sweeteners, and other fascinating\par
extracts are only a fraction of the high-profit selection. The available literature includes\par
pamphlets extolling the amazing return of youth one can anticipate while drinking a potion\par
steeped with tropical weeds, as well as volumes assuring the reader of an almost eternal lon-\par
gevity.\par
The store is directly keyed to arouse visitors' concern over their health and to capitalize on\par
real and imagined problems by offering solutions that, incidentally, cost more than the cus-\par
tomer may be able to afford. Health food store patrons are often cajoled into buying tonics\par
that promise to make the functioning of healthy organs even better, regardless of whether\par
an improvement is called for. Promotion of expensive products that consumers do not actu-\par
ally need takes considerable initiative and insight. On occasion, there may even be some\par
slight disregard for truth in an entrepreneur's zeal to cure customers of ills for a price.\par
\endgroup
\medskip
\noindent\textbf{25.} Which of the following is the main topic of the passage?
\begin{enumerate}[label=(\Alph*),leftmargin=*,itemsep=0.2em]
\item Invigorating claims regarding health food
\item Praising the health food store inventory
\item Proving the wonders of health food products
\item Marketing bogus miracles in health food stores
\end{enumerate}
\par
\begingroup
\setlength{\parskip}{0pt}
\setlength{\parindent}{0pt}
\textbf{Review}\par
Answer: D\par
Confidence: high\par
Passage ini berfokus pada pemasaran menyesatkan terkait klaim mirip mukjizat di toko makanan sehat.\par
\endgroup
\par
\vspace{0.25\baselineskip}
\noindent\textbf{26.} Which of the following best describes the author's tone?
\begin{enumerate}[label=(\Alph*),leftmargin=*,itemsep=0.2em]
\item Approving
\item Factual
\item Sarcastic
\item Hesitant
\end{enumerate}
\par
\begingroup
\setlength{\parskip}{0pt}
\setlength{\parindent}{0pt}
\textbf{Review}\par
Answer: C\par
Confidence: high\par
Penulis memakai bahasa yang skeptis dan tajam, jadi tonenya sarcastic.\par
\endgroup
\par
\vspace{0.25\baselineskip}
\noindent\textbf{27.} In line 2, the word "touted" is closest in meaning to
\begin{enumerate}[label=(\Alph*),leftmargin=*,itemsep=0.2em]
\item talked about
\item figured out
\item identified
\item known
\end{enumerate}
\par
\begingroup
\setlength{\parskip}{0pt}
\setlength{\parindent}{0pt}
\textbf{Review}\par
Answer: A\par
Confidence: medium\par
Evidence refs: line 2\par
Evidence quote: "touted"\par
Touted paling dekat maknanya dengan talked about atau promoted.\par
\endgroup
\par
\vspace{0.25\baselineskip}
\noindent\textbf{28.} Why does the author mention garlic and squash?
\begin{enumerate}[label=(\Alph*),leftmargin=*,itemsep=0.2em]
\item To explain their prevalence in diets of some ethnic groups
\item To promote their sales as healing agents for various ills
\item To compare them to modern and beneficial health products
\item To exemplify the persistence of misconceptions regarding food
\end{enumerate}
\par
\begingroup
\setlength{\parskip}{0pt}
\setlength{\parindent}{0pt}
\textbf{Review}\par
Answer: D\par
Confidence: high\par
Garlic dan squash adalah contoh miskonsepsi yang bertahan tentang kekuatan penyembuhan makanan.\par
\endgroup
\par
\vspace{0.25\baselineskip}
\noindent\textbf{29.} Where in the passage does the author give reasons for the commercial success of the health food industry?
\begin{enumerate}[label=(\Alph*),leftmargin=*,itemsep=0.2em]
\item Lines 1-2
\item Lines 8-9
\item Lines 10-13
\item Lines 14-15
\end{enumerate}
\par
\begingroup
\setlength{\parskip}{0pt}
\setlength{\parindent}{0pt}
\textbf{Review}\par
Answer: D\par
Confidence: medium\par
Evidence refs: lines 14-15\par
Evidence quote: "capitalize on real and imagined problems"\par
Kesuksesan komersial dijelaskan saat passage menerangkan bagaimana toko mengeksploitasi kekhawatiran pelanggan.\par
\endgroup
\par
\vspace{0.25\baselineskip}
\noindent\textbf{30.} In line 16, the word "cajoled" is closest in meaning to
\begin{enumerate}[label=(\Alph*),leftmargin=*,itemsep=0.2em]
\item trained
\item frightened
\item drilled
\item coaxed
\end{enumerate}
\par
\begingroup
\setlength{\parskip}{0pt}
\setlength{\parindent}{0pt}
\textbf{Review}\par
Answer: D\par
Confidence: high\par
Evidence refs: line 16\par
Evidence quote: "cajoled"\par
Cajoled berarti coaxed atau persuaded.\par
\endgroup
\par
\vspace{0.25\baselineskip}
\noindent\textbf{31.} It can be inferred from the passage that health food store operators are primarily concerned with
\begin{enumerate}[label=(\Alph*),leftmargin=*,itemsep=0.2em]
\item persuading their customers of the high quality of their wares
\item arriving at long-term solutions for health maintenance
\item maximizing profits by taking advantage of consumer naivete
\item exposing the grave consequences of neglecting one's health
\end{enumerate}
\par
\begingroup
\setlength{\parskip}{0pt}
\setlength{\parindent}{0pt}
\textbf{Review}\par
Answer: C\par
Confidence: high\par
Evidence refs: lines 15-20\par
Evidence quote: "cost more ... do not actually need"\par
Passage ini menyiratkan bahwa pemilik toko memaksimalkan keuntungan dengan mengeksploitasi gullibility konsumen.\par
\endgroup
\par
\vspace{0.25\baselineskip}
\noindent\textbf{32.} The author would most probably agree with which of the following statements?
\begin{enumerate}[label=(\Alph*),leftmargin=*,itemsep=0.2em]
\item Health food articles are positively exotic and exorbitant.
\item Promoting and selling health foods verges on cheating.
\item Health food enterprises are dedicated to absolute honesty.
\item Inducing patrons to buy health products is criminal at best.
\end{enumerate}
\par
\begingroup
\setlength{\parskip}{0pt}
\setlength{\parindent}{0pt}
\textbf{Review}\par
Answer: B\par
Confidence: high\par
Evidence refs: lines 19-20\par
Evidence quote: "slight disregard for truth"\par
Penulis menyiratkan bisnis ini bisa mendekati cheating.\par
\endgroup
\par
\vspace{0.25\baselineskip}
\subsection*{Questions 33-41}
\begingroup
\small
\sloppy
\setlength{\parskip}{0pt}
\setlength{\parindent}{0pt}
Because geologists have long indicated that fossil fuels will not last indefinitely, the U.S.\par
government finally acknowledged that sooner or later other energy sources would be needed\par
and, as a result, turned its attention to nuclear power. It was anticipated that nuclear power\par
plants could supply electricity in such large amounts and so inexpensively that they would\par
be integrated into an economy in which electricity would take over virtually all fuel-gener-\par
ating functions at nominal costs. Thus, the government subsidized the promotion of com-\par
mercial nuclear power plants and authorized their construction by utility companies. In the\par
1960s and early 1970s, the public accepted the notion of electricity being generated by nu-\par
clear reactors, and the Nuclear Regulatory Commission proceeded with plans for numerous\par
nuclear power plants in or near residential areas. By 1975, 54 plants were fully operational,\par
supplying 11 percent of the nation's electricity, and another 167 plants were at various\par
stages of planning and construction. Officials estimated that by 1990 hundreds of plants\par
would be on line, and by the turn of the century as many as 1,000 plants would be in work-\par
ing order.\par
Since 1975, this outlook and this estimation have changed drastically, and many utilities\par
have canceled existing orders. In some cases, construction was terminated even after bil-\par
lions of dollars had already been invested. After being completed and licensed at a cost of al-\par
most \$6 billion, the Shoreham Power Plant on Long Island was turned over to the state of\par
New York to be dismantled without ever having generated electric power. The reason was\par
that residents and state authorities deemed that there was no possibility of evacuating resi-\par
dents from the area should an accident occur.\par
Just 68 of those plants under way in 1975 have been completed, and another 3 are still un-\par
der construction. Therefore, it appears that in the mid 1990s, 124 nuclear power plants in\par
the nation will be in operation, generating about 18 percent of the nation's electricity, a fig-\par
ure that will undoubtedly decline as relatively outdated plants are shut down.\par
\endgroup
\medskip
\noindent\textbf{33.} What was initially planned for the nation's fuel supply in the 1950s and in the early 1960s?
\begin{enumerate}[label=(\Alph*),leftmargin=*,itemsep=0.2em]
\item Expansion and renovation of existing fuel-generating plants
\item Creation of additional storage capacities for fossil fuels
\item Conversion of the industry and the economy to nuclear power
\item Development of an array of alternative fuel and power sources
\end{enumerate}
\par
\begingroup
\setlength{\parskip}{0pt}
\setlength{\parindent}{0pt}
\textbf{Review}\par
Answer: C\par
Confidence: high\par
Evidence refs: lines 3-6\par
Evidence quote: "electricity would take over virtually all fuel-generating functions"\par
Rencananya adalah mengalihkan industri dan ekonomi secara besar-besaran ke nuclear power.\par
\endgroup
\par
\vspace{0.25\baselineskip}
\noindent\textbf{34.} How does the author describe the attitude of the population in regard to nuclear power as fuel in the early to mid 1970s?
\begin{enumerate}[label=(\Alph*),leftmargin=*,itemsep=0.2em]
\item Apprehensive
\item Ambivalent
\item Receptive
\item Resentful
\end{enumerate}
\par
\begingroup
\setlength{\parskip}{0pt}
\setlength{\parindent}{0pt}
\textbf{Review}\par
Answer: C\par
Confidence: high\par
Evidence refs: lines 8-10\par
Evidence quote: "the public accepted the notion"\par
Kata accepted menunjukkan bahwa publik bersikap receptive.\par
\endgroup
\par
\vspace{0.25\baselineskip}
\noindent\textbf{35.} In line 6, the word "nominal" is closest in meaning to
\begin{enumerate}[label=(\Alph*),leftmargin=*,itemsep=0.2em]
\item so-called
\item minimal
\item exorbitant
\item inflated
\end{enumerate}
\par
\begingroup
\setlength{\parskip}{0pt}
\setlength{\parindent}{0pt}
\textbf{Review}\par
Answer: B\par
Confidence: high\par
Evidence refs: line 6\par
Evidence quote: "nominal costs"\par
Nominal costs berarti biaya minimal.\par
\endgroup
\par
\vspace{0.25\baselineskip}
\noindent\textbf{36.} In line 8, the word "notion" is closest in meaning to
\begin{enumerate}[label=(\Alph*),leftmargin=*,itemsep=0.2em]
\item nonsense
\item notice
\item idea
\item consequence
\end{enumerate}
\par
\begingroup
\setlength{\parskip}{0pt}
\setlength{\parindent}{0pt}
\textbf{Review}\par
Answer: C\par
Confidence: high\par
Evidence refs: line 8\par
Evidence quote: "notion"\par
Notion di sini berarti idea.\par
\endgroup
\par
\vspace{0.25\baselineskip}
\noindent\textbf{37.} In line 15, the phrase "this outlook" refers to
\begin{enumerate}[label=(\Alph*),leftmargin=*,itemsep=0.2em]
\item the number of operating nuclear plants
\item the expectation for the increase in the number of nuclear plants
\item the possibility of generating electricity at nuclear installations
\item the forecast for the capacity of the nuclear plants
\end{enumerate}
\par
\begingroup
\setlength{\parskip}{0pt}
\setlength{\parindent}{0pt}
\textbf{Review}\par
Answer: B\par
Confidence: high\par
Evidence refs: line 15\par
Evidence quote: "this outlook"\par
This outlook merujuk pada ekspektasi peningkatan jumlah nuclear plants.\par
\endgroup
\par
\vspace{0.25\baselineskip}
\noindent\textbf{38.} It can be inferred from the passage that government-officials made a critical error in judgment by
\begin{enumerate}[label=(\Alph*),leftmargin=*,itemsep=0.2em]
\item disregarding the low utility of nuclear power plants
\item relying on inferior materials and faulty plant design
\item overlooking the possibility of a meltdown, however remote
\item locating installations in densely wooded areas
\end{enumerate}
\par
\begingroup
\setlength{\parskip}{0pt}
\setlength{\parindent}{0pt}
\textbf{Review}\par
Answer: C\par
Confidence: medium\par
Evidence refs: lines 19-21\par
Evidence quote: "no possibility of evacuating residents"\par
Contoh Shoreham menunjukkan bahwa pejabat meremehkan konsekuensi kecelakaan serius.\par
\endgroup
\par
\vspace{0.25\baselineskip}
\noindent\textbf{39.} The author of the passage implies that the construction of new nuclear power plants
\begin{enumerate}[label=(\Alph*),leftmargin=*,itemsep=0.2em]
\item is continuing on a smaller scale
\item is being geared for greater safety
\item has been completely halted for fear of disaster
\item has been decelerated but not terminated
\end{enumerate}
\par
\begingroup
\setlength{\parskip}{0pt}
\setlength{\parindent}{0pt}
\textbf{Review}\par
Answer: A\par
Confidence: high\par
Evidence refs: lines 22-25\par
Evidence quote: "another 3 are still under construction"\par
Passage ini menunjukkan bahwa konstruksi masih berlanjut, tetapi dalam skala yang jauh lebih kecil.\par
\endgroup
\par
\vspace{0.25\baselineskip}
\noindent\textbf{40.} Which of the following best describes the organization of the passage?
\begin{enumerate}[label=(\Alph*),leftmargin=*,itemsep=0.2em]
\item The exposition of the public opinion polls on nuclear power
\item A narration of power-source deliberation in nuclear power plants
\item Causal connections in the government's position on nuclear power
\item Point and counterpoint in the nuclear power debate
\end{enumerate}
\par
\begingroup
\setlength{\parskip}{0pt}
\setlength{\parindent}{0pt}
\textbf{Review}\par
Answer: C\par
Confidence: high\par
Passage ini disusun melalui hubungan cause-and-effect antara kebijakan dan hasilnya.\par
\endgroup
\par
\vspace{0.25\baselineskip}
\noindent\textbf{41.} The author of the passage implies that the issue of finding adequate sources of fuel and power for the future
\begin{enumerate}[label=(\Alph*),leftmargin=*,itemsep=0.2em]
\item has long been ignored by short-sighted government authorities
\item may be condoned by vacillating officials
\item has lost its pertinence in light of new discoveries
\item has not yet been satisfactorily resolved
\end{enumerate}
\par
\begingroup
\setlength{\parskip}{0pt}
\setlength{\parindent}{0pt}
\textbf{Review}\par
Answer: D\par
Confidence: high\par
Evidence refs: lines 22-25\par
Evidence quote: "the issue ... has not yet been satisfactorily resolved"\par
Masalah bahan bakar dan energi di masa depan masih belum terselesaikan.\par
\endgroup
\par
\vspace{0.25\baselineskip}
\subsection*{Questions 42-50}
\begingroup
\small
\sloppy
\setlength{\parskip}{0pt}
\setlength{\parindent}{0pt}
Collecting maps can be an enjoyable hobby for antiquarian booksellers, a captivating in-\par
terest for cartographers, a lucrative vocation for astute dealers, and an inspirational part of\par
the occupational functioning of map catalogers, archivists, and historians. Among recog-\par
nized collectibles, maps are relatively rarer than stamps, but they have had their avid enthu-\par
siasts and admirers ever since copies were made by hand only for the affluent, the command-\par
ing officer, and the ship captain.\par
Whether the interest is business-related or amateur, the economic means abundant or\par
slim, a collection needs a theme, be it associated with contemporary changes in cartographic\par
representation or geographic knowledge, or a more accessible goal centered on a particular\par
mapmaker, technique, or type of subject matter. Collectors should not overlook topical\par
maps issued predominantly or exclusively after World War II, such as navigational charts,\par
industrial compound road layouts, or aerial projections. Potential collectors ought not to\par
disregard two superficially prosaic, yet important themes: maps of travel routes for family\par
trips, and maps that, for aesthetic reasons, they personally find intriguing or simply attrac-\par
tive. In the first case, like the box with old family photos, the collection will give the travel-\par
ers the opportunity to reminisce and relive the journey.\par
In most cases, photocopies are worthy alternatives to originals. For example, historical\par
society collections customarily include the high quality facsimiles needed to make a collec-\par
tion as comprehensive and practical as possible, supplementing the contributions made by\par
well-to-do donors and benefactors. If not predisposed to wait patiently, and possibly inef-\par
fectually, for a lucky find, collectors may choose to sift through dealer stock, peruse through\par
advertisements in local, regional, or national periodicals, and solicit the assistance of the\par
U.S. Library of Congress and private agencies. Government and public agencies, compa-\par
nies, and trade associations can advise the collector about maps currently in circulation and\par
pending sales of dated reproductions, editions, and prints.\par
\endgroup
\medskip
\noindent\textbf{42.} What is the main idea of the passage?
\begin{enumerate}[label=(\Alph*),leftmargin=*,itemsep=0.2em]
\item Why hobbyists always flaunt their map collections
\item How maps can be collected by professionals and enthusiasts
\item How to assure an uninterrupted flow of collectibles
\item What cartographers advocate as a worthy undertaking
\end{enumerate}
\par
\begingroup
\setlength{\parskip}{0pt}
\setlength{\parindent}{0pt}
\textbf{Review}\par
Answer: B\par
Confidence: high\par
Passage ini menjelaskan bahwa peta dapat dikoleksi oleh profesional maupun enthusiast.\par
\endgroup
\par
\vspace{0.25\baselineskip}
\noindent\textbf{43.} In line 2, the word "lucrative" is closest in meaning to
\begin{enumerate}[label=(\Alph*),leftmargin=*,itemsep=0.2em]
\item instructive
\item insensitive
\item profitable
\item profuse
\end{enumerate}
\par
\begingroup
\setlength{\parskip}{0pt}
\setlength{\parindent}{0pt}
\textbf{Review}\par
Answer: C\par
Confidence: high\par
Evidence refs: line 2\par
Evidence quote: "lucrative"\par
Lucrative berarti profitable.\par
\endgroup
\par
\vspace{0.25\baselineskip}
\noindent\textbf{44.} According to the passage, map collecting as a hobby is
\begin{enumerate}[label=(\Alph*),leftmargin=*,itemsep=0.2em]
\item not deserving of the time and resources
\item not as conventional as collecting stamps
\item as eccentric as collecting dolls
\item conformist in the best sense of the word
\end{enumerate}
\par
\begingroup
\setlength{\parskip}{0pt}
\setlength{\parindent}{0pt}
\textbf{Review}\par
Answer: B\par
Confidence: high\par
Evidence refs: line 4\par
Evidence quote: "maps are relatively rarer than stamps"\par
Hobi ini digambarkan kurang konvensional dibanding koleksi perangko.\par
\endgroup
\par
\vspace{0.25\baselineskip}
\noindent\textbf{45.} It can be inferred from the passage that, at a time when maps were accessible to the upper socioeconomic classes, they appealed also to a fair number of
\begin{enumerate}[label=(\Alph*),leftmargin=*,itemsep=0.2em]
\item professional copiers
\item ardent devotees
\item buried-treasure hunters
\item obscure amateur dealers
\end{enumerate}
\par
\begingroup
\setlength{\parskip}{0pt}
\setlength{\parindent}{0pt}
\textbf{Review}\par
Answer: B\par
Confidence: medium\par
Evidence refs: lines 4-6\par
Evidence quote: "avid enthusiasts and admirers"\par
Passage ini menyiratkan bahwa peta juga menarik bagi ardent devotees.\par
\endgroup
\par
\vspace{0.25\baselineskip}
\noindent\textbf{46.} In line 7, the phrase "economic means" is closest in meaning to
\begin{enumerate}[label=(\Alph*),leftmargin=*,itemsep=0.2em]
\item economic maps
\item fiscal responsibility
\item available funds
\item capital investment
\end{enumerate}
\par
\begingroup
\setlength{\parskip}{0pt}
\setlength{\parindent}{0pt}
\textbf{Review}\par
Answer: C\par
Confidence: high\par
Evidence refs: line 7\par
Evidence quote: "economic means"\par
Economic means merujuk pada available funds.\par
\endgroup
\par
\vspace{0.25\baselineskip}
\noindent\textbf{47.} The author of the passage mentions all of the following as sources of procuring maps EXCEPT:
\begin{enumerate}[label=(\Alph*),leftmargin=*,itemsep=0.2em]
\item fellow collectors
\item map vendors
\item personal archives
\item publishers
\end{enumerate}
\par
\begingroup
\setlength{\parskip}{0pt}
\setlength{\parindent}{0pt}
\textbf{Review}\par
Answer: A\par
Confidence: medium\par
Evidence refs: lines 21-24\par
Evidence quote: "dealer stock ... advertisements ... Library of Congress and private agencies"\par
Passage ini menyebut beberapa sumber, tetapi tidak menyebut fellow collectors.\par
\endgroup
\par
\vspace{0.25\baselineskip}
\noindent\textbf{48.} In line 13, the author uses the phrase "superficially prosaic" to mean
\begin{enumerate}[label=(\Alph*),leftmargin=*,itemsep=0.2em]
\item described in informal prose
\item seemingly boring and unimaginative
\item useful for travelers who enjoy a change
\item potentially uncovered in a box of photos
\end{enumerate}
\par
\begingroup
\setlength{\parskip}{0pt}
\setlength{\parindent}{0pt}
\textbf{Review}\par
Answer: B\par
Confidence: high\par
Evidence refs: line 13\par
Evidence quote: "superficially prosaic"\par
Frasa ini berarti tampak membosankan dan tidak imajinatif.\par
\endgroup
\par
\vspace{0.25\baselineskip}
\noindent\textbf{49.} In line 20, the word "predisposed" is closest in meaning to
\begin{enumerate}[label=(\Alph*),leftmargin=*,itemsep=0.2em]
\item pressured
\item provoked
\item condemned
\item inclined
\end{enumerate}
\par
\begingroup
\setlength{\parskip}{0pt}
\setlength{\parindent}{0pt}
\textbf{Review}\par
Answer: D\par
Confidence: high\par
Evidence refs: line 20\par
Evidence quote: "predisposed"\par
Predisposed paling dekat maknanya dengan inclined.\par
\endgroup
\par
\vspace{0.25\baselineskip}
\noindent\textbf{50.} A paragraph following the passage would most likely discuss
\begin{enumerate}[label=(\Alph*),leftmargin=*,itemsep=0.2em]
\item specific organizations to contact about map acquisition
\item specific mapping techniques used to enlarge the scale
\item trimming and framing valuable acquisitions
\item volunteering time and work to maintain obsolete maps
\end{enumerate}
\par
\begingroup
\setlength{\parskip}{0pt}
\setlength{\parindent}{0pt}
\textbf{Review}\par
Answer: A\par
Confidence: high\par
Evidence refs: lines 21-25\par
Evidence quote: "Library of Congress and private agencies"\par
Paragraf berikutnya kemungkinan melanjutkan dengan organisasi spesifik untuk perolehan peta.\par
\endgroup
\par
\vspace{0.25\baselineskip}
\newpage
\section*{Answer Key Appendix}
\subsection*{Listening}
\begin{itemize}[leftmargin=*]
\item 1: A
\item 2: C
\item 3: B
\item 4: A
\item 5: B
\item 6: C
\item 7: A
\item 8: B
\item 9: D
\item 10: D
\item 11: C
\item 12: A
\item 13: A
\item 14: D
\item 15: B
\item 16: A
\item 17: A
\item 18: C
\item 19: B
\item 20: D
\item 21: A
\item 22: B
\item 23: C
\item 24: C
\item 25: B
\item 26: C
\item 27: D
\item 28: B
\item 29: D
\item 30: B
\item 31: B
\item 32: D
\item 33: B
\item 34: B
\item 35: B
\item 36: B
\item 37: D
\item 38: C
\item 39: D
\item 40: C
\item 41: C
\item 42: B
\item 43: D
\item 44: A
\item 45: B
\item 46: A
\item 47: B
\item 48: B
\item 49: B
\item 50: A
\end{itemize}
\subsection*{Structure}
\begin{itemize}[leftmargin=*]
\item 1: B
\item 2: A
\item 3: C
\item 4: C
\item 5: B
\item 6: D
\item 7: B
\item 8: A
\item 9: B
\item 10: B
\item 11: C
\item 12: B
\item 13: C
\item 14: D
\item 15: B
\item 16: D
\item 17: D
\item 18: A
\item 19: D
\item 20: D
\item 21: B
\item 22: C
\item 23: B
\item 24: B
\item 25: C
\item 26: D
\item 27: D
\item 28: B
\item 29: D
\item 30: B
\item 31: D
\item 32: B
\item 33: C
\item 34: B
\item 35: B
\item 36: D
\item 37: C
\item 38: D
\item 39: B
\item 40: C
\end{itemize}
\subsection*{Reading}
\begin{itemize}[leftmargin=*]
\item 1: B
\item 2: C
\item 3: C
\item 4: A
\item 5: C
\item 6: C
\item 7: B
\item 8: B
\item 9: C
\item 10: D
\item 11: B
\item 12: B
\item 13: B
\item 14: A
\item 15: C
\item 16: B
\item 17: C
\item 18: B
\item 19: B
\item 20: C
\item 21: B
\item 22: C
\item 23: D
\item 24: B
\item 25: D
\item 26: C
\item 27: A
\item 28: D
\item 29: D
\item 30: D
\item 31: C
\item 32: B
\item 33: C
\item 34: C
\item 35: B
\item 36: C
\item 37: B
\item 38: C
\item 39: A
\item 40: C
\item 41: D
\item 42: B
\item 43: C
\item 44: B
\item 45: B
\item 46: C
\item 47: A
\item 48: B
\item 49: D
\item 50: A
\end{itemize}
\end{document}
